\subsection{The Energy-Conjugate Trace}
\label{sec:trace-energy}

The energy-conjugate trace is defined by requiring that the beam strains 
$\bm{e} = (\epsilon, \bm{\gamma}_\perp, \varkappa, \bm{\kappa}_\perp)$ 
be work-conjugate to the stress resultants 
$\bm{s} = (N, \bm{F}, T, \bm{M})$ \eqref{eq:energy-conjugacy-def}:
\begin{equation*}
\delta\bm{e} \cdot \snm = \int_\Section \delta\bm{\varepsilon} \cdot \bm{\sigma} \dA
\end{equation*}
for all variations within the Saint-Venant class. 
The stress resultants are:
\begin{equation*}
N := \int_\Section \sigma \dA, \qquad 
\bm{F} := \int_\Section \bm{\tau} \dA, \qquad 
T := \int_\Section (\FrameBasis \times \bm{r}) \cdot \bm{\tau} \dA, \qquad 
\bm{M} := \int_\Section \bm{r} \times \bm{\sigma} \dA,
% \label{eq:resultant-defs}
\end{equation*}
where $\bm{\sigma} = \sigma\FrameBasis + \bm{\tau}$ is the traction on the section.

\iffalse
\subsubsection{Resultant Evolution}

The stress resultants for a Saint-Venant solution evolve affinely along the beam:
\begin{equation*}
\bm{s}(x) = \mathbf{R}(x)\,\bm{p} = (\IesanStiff + x\mathbf{R}_1)\,\bm{p}.
\end{equation*}
where \(\IesanStiff\) is defined by \eqref{eq:R0-matrix}. 
The constant part encodes the constitutive relations at $x = 0$.
The $x$-linear part arises from moment equilibrium $\bm{M}' + \FrameBasis \times \bm{F}=\mathbf{0}$:
\begin{equation*}
\mathbf{R}_1 = \mathbb{J}\IesanStiff =\begin{pmatrix}
0 & \mathbf{0}^\top & 0 & \mathbf{0}^\top \\[4pt]
\mathbf{0} & \mathbf{0} & \mathbf{0} & \mathbf{0} \\[4pt]
0 & \mathbf{0}^\top & 0 & \mathbf{0}^\top \\[4pt]
\mathbf{0} & \mathbf{0} & \mathbf{0} & -E\FrameBasis^\times\IntRGoRG
\end{pmatrix}.
% \label{eq:R1-matrix}
\end{equation*}
\fi 


% 
\subsubsection{Strain Energy}

The strain energy at section $x$ is quadratic in the evolved parameters 
$\exp({x\IesanEvolveSP})\bm{p}$:
\begin{equation*}
\mathcal{U}(x) = \frac{1}{2} \int_\Section \bm{\varepsilon}\cdot \mathbb{C}\, \bm{\varepsilon} \dA 
= \frac{1}{2} \bm{p}\cdot \left(\mathbf{1} + x\mathbf{G}_{\mathcal{T}^\top}\right) \IesanEnergy\, \left(\mathbf{1} + x\mathbf{A}_T\right) \bm{p},
\end{equation*}
where the Saint-Venant energy matrix is
\begin{equation}
\begin{aligned}
\IesanEnergy := \int_\Section \mathbf{S}_0^\top \mathbb{C}\, \mathbf{S}_0 \dA 
% &= \begin{pmatrix}
% EA & \mathbf{0}^\top & 0 & EA\CentroidLocation^\top \\[4pt]
% \mathbf{0} & G\mathbf{H}_{bb} & G\mathbf{h}_{\varphi b} & \mathbf{0} \\[4pt]
% 0 & G\mathbf{h}_{\varphi b}^\top & G\Jsv & \mathbf{0}^\top \\[4pt]
% EA\CentroidLocation & \mathbf{0} & \mathbf{0} & E\IntRoR \\
% \end{pmatrix}\\
% &= \begin{pmatrix}
% EA & EA\CentroidLocation^\top & \mathbf{0}^\top & 0 \\[4pt]
% EA\CentroidLocation & E\IntRoR & \mathbf{0} & \mathbf{0} \\
% 0  & \mathbf{0}^\top & G\Jsv & G\mathbf{h}_{\varphi b}^\top \\[4pt]
% \mathbf{0}& \mathbf{0}& G\mathbf{h}_{\varphi b} & G\mathbf{H}_{bb}   \\[4pt]
% \end{pmatrix} \\
&=\begin{pmatrix}
EA & EA\CentroidLocation^\top & 0 & \mathbf{0}^\top \\
EA\CentroidLocation & E\IntRoR & \mathbf{0} & \mathbf{0} \\
0 & \mathbf{0}^\top & G\Jsv & G(\Jsv\CenterShearLocation + \mathbf{h}_{\varphi b})^\top \\
\mathbf{0} & \mathbf{0} & G(\Jsv\CenterShearLocation + \mathbf{h}_{\varphi b}) & G(\mathbf{H}_{bb} + \Jsv\CenterShearLocation\otimes\CenterShearLocation + \CenterShearLocation\otimes\mathbf{h}_{\varphi b} + \mathbf{h}_{\varphi b}\otimes\CenterShearLocation)
\end{pmatrix}
\end{aligned}
\label{eq:energy-matrix}
\end{equation}
%
where shear-torsion coupling terms are:
\begin{align}
\mathbf{H}_{bb} &:= \int_\Section \mathbf{B}_{(b)}^\top \mathbf{B}_{(b)} \dA, 
\label{eq:def-H-bb} \\[4pt]
\mathbf{h}_{\varphi b} &:= \int_\Section \mathbf{B}_{(b)}^\top \mathbf{b}_{(\varphi)} \dA 
= -\Jsv\CenterShearLocation + \tfrac{\nu}{2}\bm{\zeta} - \nu\bm{\lambda},
\label{eq:h-coupling}
\end{align}
and $\bm{\zeta}$ and $\bm{\lambda}$ capture Poisson-torsion interactions:
\begin{equation}
\bm{\zeta} := \int_\Section r^2 \nabla\WarpTwistIesan \dA, \qquad 
\bm{\lambda} := \int_\Section (\bm{r} - \CentroidLocation)(\bm{r} \cdot \nabla\WarpTwistIesan) \dA.
\label{eq:def-zeta-lambda}
\end{equation}

\subsubsection{Trace Matrix}

The inverse of the energy matrix \eqref{eq:energy-matrix} is:
\begin{equation*}
\begin{aligned}
\IesanEnergy^{-1}
% &= \begin{pmatrix}
% \frac{1}{EA} + \CentroidLocation\cdot(E\IntRGoRG)^{-1}\CentroidLocation 
% & \mathbf{0}^\top & 0 & -\CentroidLocation^\top(E\IntRGoRG)^{-1} \\[6pt]
% \mathbf{0} & \frac{1}{G}\big(\mathbf{H}_{bb}^{-1} + \frac{1}{\tilde{\Jsv}}\bm{\eta}\otimes\bm{\eta}\big) 
% & -\frac{1}{G\tilde{\Jsv}}\bm{\eta} & \mathbf{0} \\[6pt]
% 0 & -\frac{1}{G\tilde{\Jsv}}\bm{\eta}^\top & \frac{1}{G\tilde{\Jsv}} & \mathbf{0}^\top \\[6pt]
% -(E\IntRGoRG)^{-1}\CentroidLocation & \mathbf{0} & \mathbf{0} & (E\IntRGoRG)^{-1}
% \end{pmatrix}\\
% &=\begin{pmatrix}
% \frac{1}{EA} + \CentroidLocation\cdot(E\IntRGoRG)^{-1}\CentroidLocation 
% & \mathbf{0}^\top & 0 & -\CentroidLocation^\top(E\IntRGoRG)^{-1} \\[6pt]
% -(E\IntRGoRG)^{-1}\CentroidLocation & \mathbf{0} & \mathbf{0} & (E\IntRGoRG)^{-1} \\[6pt]
% 0 & -\frac{1}{G\tilde{\Jsv}}\bm{\eta}^\top & \frac{1}{G\tilde{\Jsv}} & \mathbf{0}^\top \\[6pt]
% \mathbf{0} & \frac{1}{G}\big(\mathbf{H}_{bb}^{-1} + \frac{1}{\tilde{\Jsv}}\bm{\eta}\otimes\bm{\eta}\big) 
% & -\frac{1}{G\tilde{\Jsv}}\bm{\eta} & \mathbf{0} \\[6pt]
% \end{pmatrix}\\
%
% &=\begin{pmatrix}
% \frac{1}{EA} + \CentroidLocation\cdot(E\IntRGoRG)^{-1}\CentroidLocation 
% & -\CentroidLocation^\top(E\IntRGoRG)^{-1}  & 0 & \mathbf{0}^\top \\[6pt]
% -(E\IntRGoRG)^{-1}\CentroidLocation & (E\IntRGoRG)^{-1} & \mathbf{0} & \mathbf{0}  \\[6pt]
% 0  & \mathbf{0}^\top & \frac{1}{G\tilde{\Jsv}} & -\frac{1}{G\tilde{\Jsv}}\bm{\eta}^\top  \\[6pt]
% \mathbf{0} & \mathbf{0}& -\frac{1}{G\tilde{\Jsv}}\bm{\eta} & \frac{1}{G}\big(\mathbf{H}_{bb}^{-1} + \frac{1}{\tilde{\Jsv}}\bm{\eta}\otimes\bm{\eta}\big) 
%  \\[6pt]
% \end{pmatrix}\\
&=\begin{pmatrix}
\frac{1}{EA} + \CentroidLocation\cdot (E\IntRGoRG)^{-1}\CentroidLocation & -\CentroidLocation^\top (E\IntRGoRG)^{-1} & 0 & \mathbf{0}^\top \\[6pt]
-(E\IntRGoRG)^{-1}\CentroidLocation & (E\IntRGoRG)^{-1} & \mathbf{0} & \mathbf{0} \\[6pt]
0 & \mathbf{0}^\top & \frac{1}{G\Jsv} + \bm{\eta}\cdot(G\tilde{\mathbf{H}}_{bb})^{-1}\bm{\eta} 
& -\bm{\eta}^\top(G\tilde{\mathbf{H}}_{bb})^{-1} \\[6pt]
\mathbf{0} & \mathbf{0} & -(G\tilde{\mathbf{H}}_{bb})^{-1}\bm{\eta} & (G\tilde{\mathbf{H}}_{bb})^{-1}
\end{pmatrix} 
\\
&= \begin{pmatrix}
\frac{1}{EA}\FrameBasis\otimes\FrameBasis - \FrameBasis\otimes(E\IntRGoRG)^{-1}\CentroidLocation-(E\IntRGoRG)^{-1}\CentroidLocation\otimes\FrameBasis + (E\IntRGoRG)^{-1} & \mathbf{0}\\
    \mathbf{0} & 
\end{pmatrix}
\end{aligned}
% \label{eq:Cbar-inv}
\end{equation*}
where $\tilde{\mathbf{H}}_{bb} := \mathbf{H}_{bb} - \frac{1}{\Jsv}\mathbf{h}_{\varphi b}\otimes\mathbf{h}_{\varphi b}$ 
and
\begin{equation}
\tilde{\Jsv} := \Jsv - \mathbf{h}_{\varphi b}\cdot \mathbf{H}_{bb}^{-1} \mathbf{h}_{\varphi b}, 
\qquad 
\bm{\eta} := \mathbf{H}_{bb}^{-1} \mathbf{h}_{\varphi b}.
\label{eq:reduced-torsion}
\end{equation}

Using \(\IesanStiff\) given by \eqref{eq:G0-neg-m} in \eqref{eq:T0-formula}:
\[
\TraceDpDeRoot = \IesanEnergy^{-1}\IesanStiff^\top
=
\left(\begin{array}{cc|cc}
1 
& \mathbf{0}^\top
& 0 
& -EA\Bigl(\dfrac{1}{EA} + \CentroidLocation\cdot (E\IntRGoRG)^{-1}\CentroidLocation\Bigr)
(\FrameBasis\times\CentroidLocation)^\top
+ \CentroidLocation^\top (E\IntRGoRG)^{-1}(E\FrameBasis^\times\IntRoR)^\top
\\[8pt]
\mathbf{0} & \mathbf{0} & \mathbf{0} 
& \FrameBasis^{\times} 
%EA (E\IntRGoRG)^{-1}\CentroidLocation\otimes (\FrameBasis\times\CentroidLocation) - (E\IntRGoRG)^{-1}(E\FrameBasis^\times\IntRoR)^\top
\\[8pt] \hline 
0 
&  -\bm{\eta}^\top \ShiftB %(G\tilde{\mathbf{H}}_{bb})^{-1}(E\IntRGoRG)^\top 
& \mu_{(b)}
& \mathbf{0}^\top
\\[8pt]
  \mathbf{0} 
& \ShiftB
& \ShiftK 
& \mathbf{0}
\end{array}\right)
\]
where:
\begin{equation}
\begin{aligned}
\ShiftB &= (G\tilde{\mathbf{H}}_{bb})^{-1}(E\IntRGoRG)^\top \\
\mu_{(b)} &= 1 + \bm{\eta}\cdot G\Jsv (G\tilde{\mathbf{H}}_{bb})^{-1}\bm{\eta} \\
\ShiftK &= - G\Jsv (G\tilde{\mathbf{H}}_{bb})^{-1}\bm{\eta} \\
\end{aligned}
\end{equation}
where \(\lambda = 1 + EA\,\CentroidLocation\cdot (E\IntRGoRG)^{-1}\CentroidLocation\).


% \clearpage 
\subsubsection{Strain Energy}


\begin{equation*}
\begin{aligned}
\mathbf{S}_0(\bm{r}) &= \begin{bmatrix}
1 & \bm{r}^\top & 0 & \mathbf{0}^\top \\[4pt]
\mathbf{0} & \mathbf{0} & \mathbf{b}_{(\WarpTwistIesan)}(\bm{r}) & 
\mathbf{b}_{(\WarpTwistIesan)}(\bm{r}) \otimes \CenterShearLocation + \mathbf{B}_{(\WarpShearIesan)}(\bm{r})
\end{bmatrix} \\
&= \begin{bmatrix}
1 & \bm{r}^\top & 0 & \mathbf{0}^\top \\[4pt]
\mathbf{0} & \mathbf{0} & \FrameBasis\times\bm{r} & 
\mathbf{0} %\mathbf{b}_{(\WarpTwistIesan)}(\bm{r}) \otimes \CenterShearLocation (\bm{r})
\end{bmatrix} 
+\begin{bmatrix}
0 & \mathbf{0}^\top & 0 & \mathbf{0}^\top \\[4pt]
\mathbf{0} & \mathbf{0} & \nabla\WarpTwistIesan & 
\mathbf{b}_{(\WarpTwistIesan)}(\bm{r}) \otimes \CenterShearLocation + \mathbf{B}_{(\WarpShearIesan)}(\bm{r})
\end{bmatrix}
\end{aligned}
% \label{eq:S0-explicit}
\end{equation*}
and
\begin{equation*}
\mathbf{S}_1 = \begin{bmatrix}
    0 & \mathbf{0} & (\bm{r} - \CentroidLocation)^\top \\ 
    \mathbf{0} & \mathbf{0} & \mathbf{0}
\end{bmatrix}
\end{equation*}
where the shear strain shape functions are
\begin{subequations}
\begin{align}
\mathbf{b}_{(\WarpTwistIesan)}(\bm{r}) &\coloneq \FrameBasis \times \bm{r} + \nabla\WarpTwistIesan(\bm{r}), 
\label{eq:tau-phi} \\
\mathbf{B}_{(\WarpShearIesan)}(\bm{r}) &\coloneq (\nabla\bm{\WarpShearIesan})^{\top} + \PoissonFlexure(\bm{r}).
\label{eq:tau-psi}
\end{align}
\end{subequations}

The strain energy at section $x$ is quadratic in the evolved parameters 
$\exp({x\IesanEvolveSP})\bm{p}$:
\begin{equation*}
\mathcal{U}(x) = \frac{1}{2} \int_\Section \bm{\varepsilon}\cdot \mathbb{C}\, \bm{\varepsilon} \dA 
= \frac{1}{2} \bm{p}\cdot \left(\mathbf{1} + x\mathbf{G}_{\mathcal{T}}^\top\right) \IesanEnergy\, \left(\mathbf{1} + x\mathbf{A}_T\right) \bm{p},
\end{equation*}
where the Saint~Venant energy matrix is
\begin{equation}
\begin{aligned}
\IesanEnergy := \int_\Section \mathbf{S}_0^\top \mathbb{C}\, \mathbf{S}_0 \dA 
&=\begin{bmatrix}
EA & EA\CentroidLocation^\top & 0 & \mathbf{0} \\
EA\CentroidLocation & E\IntRoR & \mathbf{0} & \mathbf{0} \\
0 & \mathbf{0} & G\Jsv & G(\Jsv\CenterShearLocation + \mathbf{h}_{\WarpTwistIesan})^\top \\
\mathbf{0} & \mathbf{0} & G(\Jsv\CenterShearLocation + \mathbf{h}_{\WarpTwistIesan}) & G(\mathbf{H}_{\WarpShearIesan} + \Jsv\CenterShearLocation\otimes\CenterShearLocation + \CenterShearLocation\otimes\mathbf{h}_{\WarpTwistIesan} + \mathbf{h}_{\WarpTwistIesan}\otimes\CenterShearLocation)
\end{bmatrix}
\end{aligned}
\label{eq:energy-matrix}
\end{equation}
%
where shear-torsion coupling terms are:
\begin{subequations}
\begin{align}
\mathbf{H}_{\WarpShearIesan} &:= \int_\Section \mathbf{B}_{(b)}^\top \mathbf{B}_{(b)} \dA, 
\label{eq:def-H-bb} \\[4pt]
\mathbf{h}_{\WarpTwistIesan} &:= \int_\Section \mathbf{B}_{(b)}^\top \mathbf{b}_{(\varphi)} \dA 
% = -\Jsv\CenterShearLocation + \tfrac{\nu}{2}\bm{\zeta} - \nu\bm{\lambda},
\label{eq:h-coupling}
\end{align}
\end{subequations}
with \(\mathbf{B}_{(b)}\) and \(\mathbf{b}_{(\varphi)}\) give by \eqref{eq:tau-psi} and \eqref{eq:tau-phi}, respectively.
\begin{Details}
, 
and $\bm{\zeta}$ and $\bm{\lambda}$ capture Poisson-torsion interactions:
\begin{equation}
\bm{\zeta} := \int_\Section r^2 \nabla\WarpTwistIesan \dA, \qquad 
\bm{\lambda} := \int_\Section (\bm{r} - \CentroidLocation)(\bm{r} \cdot \nabla\WarpTwistIesan) \dA.
\label{eq:def-zeta-lambda}
\end{equation}
\end{Details}

\subsubsection{Trace Matrix}

The inverse of the energy matrix \eqref{eq:energy-matrix} is:
\begin{equation*}
\begin{aligned}
\IesanEnergy^{-1}
&=\begin{bmatrix}
\frac{1}{EA} + \CentroidLocation\cdot (E\IntRGoRG)^{-1}\CentroidLocation 
& -\CentroidLocation^\top (E\IntRGoRG)^{-1} 
& 0 
& \mathbf{0} \\[6pt]
-(E\IntRGoRG)^{-1}\CentroidLocation 
& (E\IntRGoRG)^{-1} 
& \mathbf{0} 
& \mathbf{0} \\[6pt]
0 
& \mathbf{0} 
& % \displaystyle 
  (G\Jsv)^{-1} 
  + \TraceTF\cdot
    (G\tilde{\mathbf{H}}_{bb})^{-1}
    \TraceTF
& \displaystyle -\TraceTF^\top
    (G\tilde{\mathbf{H}}_{bb})^{-1} \\[6pt]
\mathbf{0} 
& \mathbf{0} 
& -(G\tilde{\mathbf{H}}_{bb})^{-1}\TraceTF
& (G\tilde{\mathbf{H}}_{bb})^{-1}
\end{bmatrix} 
\end{aligned}
% \label{eq:Cbar-inv}
\end{equation*}
where 
\begin{equation}
% \tilde{\Jsv} := \Jsv - \mathbf{h}_{\WarpTwistIesan}\cdot \mathbf{H}_{\WarpShearIesan}^{-1} \mathbf{h}_{\WarpTwistIesan}, 
\tilde{\mathbf{H}}_{bb} := \mathbf{H}_{\WarpShearIesan} - \frac{1}{\Jsv}\mathbf{h}_{\WarpTwistIesan}\otimes\mathbf{h}_{\WarpTwistIesan},
\qquad\text{and}\qquad
\TraceTF := \frac{1}{\Jsv}\bigl(\Jsv\CenterShearLocation + \mathbf{h}_{\WarpTwistIesan}\bigr)
\qquad 
% \bm{\eta} := \mathbf{H}_{\WarpShearIesan}^{-1} \mathbf{h}_{\WarpTwistIesan}.
\label{eq:reduced-torsion}
\end{equation}

Using \(\IesanStiff\) given by \eqref{eq:G0-neg-m} in \eqref{eq:T0-formula} furnishes the inverse trace matrix with the structure \eqref{eq:T0-inv}:
\[
\TraceDpDeRoot = \IesanEnergy^{-1}\IesanStiff^\top
=
\left[\begin{array}{cc|cc}
1 
& \mathbf{0}
& 0 
& \ShiftNK
% -EA\Bigl(\dfrac{1}{EA} + \CentroidLocation\cdot (E\IntRGoRG)^{-1}\CentroidLocation\Bigr)
% (\FrameBasis\times\CentroidLocation)^\top
% + \CentroidLocation^\top (E\IntRGoRG)^{-1}(E\FrameBasis^\times\IntRoR)^\top
\\[8pt]
\mathbf{0} & \mathbf{0} & \mathbf{0} 
& \FrameBasis^{\times} 
\\ \hline && \\[-8pt]
0 
&  \ShiftT^\top 
& \ShiftScaleT
& \mathbf{0}
\\[8pt]
  \mathbf{0} 
& \ShiftB
& \ShiftK 
& \mathbf{0}
\end{array}\right]
\]
where 
\begin{equation}
\begin{aligned}
\ShiftB &= (G\tilde{\mathbf{H}}_{bb})^{-1}(E\IntRGoRG)^\top &&\sim L^{-2} \\ 
\ShiftA &= \mathbf{0},\\
\TraceTF &= \frac{1}{\Jsv}(\Jsv \CenterShearLocation + \mathbf{h}_{\WarpTwistIesan}) &&\sim L \\
\TraceFT 
%&= G(E\IntRGoRG)^{-1}(\Jsv \CenterShearLocation + \mathbf{h}_{\WarpTwistIesan}) \sim L \\
 &= G\Jsv (E\IntRGoRG)^{-1}\TraceTF &&\sim L \\
\ShiftN &= \mathbf{0},\\[2pt]
\ShiftNK &= -\Bigl(1 + EA\,\CentroidLocation\cdot (E\IntRGoRG)^{-1}\CentroidLocation\Bigr)
  (\FrameBasis\times\CentroidLocation)
  + \FrameBasis\times E\IntRoR\,(E\IntRGoRG)^{-1}\CentroidLocation && \sim L \\
\ShiftF &= \mathbf{0},\\ 
\hline \\[-8pt]
\ShiftT 
%&= -\frac{1}{\Jsv}\,\ShiftB^\top\bigl(\Jsv\CenterShearLocation + \mathbf{h}_{\WarpTwistIesan}\bigr) \\
 &= -\ShiftB^\top\TraceTF && \sim L^{-1} \\
\ShiftK &= -\tilde{\mathbf{H}}_{bb}^{-1}\bigl(\Jsv\CenterShearLocation + \mathbf{h}_{\WarpTwistIesan}\bigr) \\
\ShiftScaleT 
% &= 1 + \frac{1}{\Jsv}\bigl(\Jsv\CenterShearLocation + \mathbf{h}_{\WarpTwistIesan}\bigr)\cdot\tilde{\mathbf{H}}_{bb}^{-1}\bigl(\Jsv\CenterShearLocation + \mathbf{h}_{\WarpTwistIesan}\bigr) \\
&= 1 - \TraceTF\cdot \ShiftK && \sim 1 \\
\ShiftScaleN &= 0
\end{aligned}
\label{eq:energetic-trace-blocks}
\end{equation}


\subsubsection{Summary}

\begin{table}[ht]
\centering
\caption{Symbols introduced in the energy-conjugate trace derivation.}
\label{tab:energy-trace-symbols}
\begin{tabular}{@{}lll@{}}
\toprule
Symbol & Units & Description \\
\midrule
\multicolumn{3}{@{}l}{\textit{Shear strain shape functions}} \\[2pt]
$\mathbf{b}_{(\varphi)}$ & L & Torsion shear shape,  \eqref{eq:tau-phi} \\
$\mathbf{B}_{(b)}$ & L & Flexural shear shape, \eqref{eq:tau-psi} \\[6pt]
\multicolumn{3}{@{}l}{\textit{Section stiffness quantities}} \\[2pt]
$\mathbf{H}_{bb}$ & \Length{6} %\Length{4} 
& Flexural shear stiffness, 
    \eqref{eq:def-H-bb} \\
$\mathbf{h}_{\varphi b}$ & \Length{5} %\Length{4}
& Shear-torsion coupling 
    \eqref{eq:h-coupling} \\
% $\tilde{\Jo}$ & \Length{4} & Reduced torsion constant \eqref{eq:reduced-torsion} \\
% $\bm{\eta}$ & \Length{-1} %--- 
% & Shear-torsion coupling ratio \eqref{eq:reduced-torsion} \\[6pt]
\multicolumn{3}{@{}l}{\textit{Poisson-torsion coupling}} \\[2pt]
% $\bm{\zeta}$ & \Length{5} & Torsion gradient moment \eqref{eq:def-zeta-lambda} \\
% $\bm{\lambda}$ & \Length{5} & Torsion-centroid coupling \eqref{eq:def-zeta-lambda} \\[6pt]
% \multicolumn{3}{@{}l}{\textit{Resultant and energy matrices}} \\[2pt]
% $\IesanEnergy$ & mixed & Canonical energy matrix \\
% $\IesanStiff$ & mixed & Stress resultant matrix at $x = 0$ \eqref{eq:G0-neg-m} \\
\bottomrule
\end{tabular}
\end{table}

\begin{Details}
\begin{proposition}[Consistency of Energy Conjugacy]
The consistency condition \eqref{eq:consistency-condition} is satisfied by 
the Saint-Venant structure.
\end{proposition}

\begin{proof}
The matrix $\mathbf{R}_1$ has only the $(4,4)$ block nonzero, corresponding 
to moment evolution under shear force: $\bm{M}' = \FrameBasis \times \bm{F}$.
Thus $\IesanStiff^{-1}\mathbf{R}_1$ acts only on the $\bm{b}_{\Section}$ columns.

A moment increment $\delta\bm{M} = -E\FrameBasis^\times\IntRGoRG\,\bm{b}_{\Section}$ 
with $\delta N = \delta\bm{F} = \delta T = 0$ must arise from Ieşan parameter 
increments. 
From $\delta N = 0$:
\begin{equation*}
0 = EA\,\delta a_\varepsilon + EA\CentroidLocation \cdot \delta\bm{a}_\Omega
\qquad \Longrightarrow \qquad
\delta a_\varepsilon = -\CentroidLocation \cdot \delta\bm{a}_\Omega.
\end{equation*}
Substituting into the moment relation furnishes:
\begin{equation*}
\delta\bm{M} = EA\FrameBasis\times\CentroidLocation(\CentroidLocation \cdot \delta\bm{a}_\Omega) 
- E\FrameBasis\times\IntRoR\,\delta\bm{a}_\Omega
= -E\FrameBasis\times\IntRGoRG\,\delta\bm{a}_\Omega.
\end{equation*}
Thus $\delta\bm{a}_\Omega = \bm{b}_{\Section}$ and 
$\delta a_\varepsilon = -\CentroidLocation \cdot \bm{b}_{\Section}$, confirming:
\begin{equation*}
\IesanStiff^{-1}\mathbf{R}_1 = \begin{pmatrix}
0 & \mathbf{0}^\top & 0 & -\CentroidLocation^\top \\
\mathbf{0} & \mathbf{0} & \mathbf{0} & \mathbf{0} \\
0 & \mathbf{0}^\top & 0 & \mathbf{0}^\top \\
\mathbf{0} & \mathbf{0} & \mathbf{0} & \mathbf{1}_\Section
\end{pmatrix} = \IesanEvolveSP. \qedhere
\end{equation*}
\end{proof}
\end{Details}

% \subsubsection{The Energy-Conjugate (Renton-Romano) Trace}

The energy-conjugate trace is defined by requiring that the beam strains 
$(\bm{\gamma}, \bm{\kappa})$ be work-conjugate to the classical resultants 
$(\bm{n}, \bm{m})$. 
That is, for any variation $\delta\hat{\bm{u}}$ within the 
Saint-Venant class:
\begin{equation}
\delta\bm{\gamma} \cdot \bm{n} + \delta\bm{\kappa} \cdot \bm{m} 
= \int_\Section \delta\bm{\varepsilon} : \bm{\sigma} \, \mathrm{d}A.
\label{eq:energy-conjugacy}
\end{equation}

We derive this trace by examining the structure of the Saint-Venant strain 
field and identifying which strain measures satisfy \eqref{eq:energy-conjugacy}.

\paragraph{Saint-Venant strain fields.}

From \eqref{eq:iesan-point-strain}, the pointwise strains in a Saint-Venant 
solution are:
\begin{equation}
    \bm{\varepsilon} = \underbrace{\begin{bmatrix}
        \FrameBasis\otimes\FrameBasis & \FrameBasis\otimes\bm{r} & \FrameBasis\times\bm{r} + \nabla\WarpTwistIesan
        & \mathbf{B}_{(b)}
    \end{bmatrix}}_{:=\mathbf{B}}
    \begin{pmatrix}
        \epsilon_0 \\ \bm{\epsilon}_1\\ \varkappa \\ \bm{\beta}
    \end{pmatrix}
\label{eq:sv-shear-strain}
\end{equation}
where the local parameters are 
$\hat{\bm{e}} = (\epsilon_0, \bm{\beta}, \hat{\varkappa}, \bm{\epsilon}_1) = \mathbf{X}^{-1}\,\bm{e}$ and 
\begin{equation*}
\mathbf{B}_{b} = (\nabla\WarpShearIesan)^T + \bm{W}_{(b)}^T
\qquad\text{with}\qquad
\bm{W}_{(b)} := -\nu\big(\mathrm{dev}(\bm{r} \otimes \bm{r}) 
- \bm{r} \otimes \CentroidLocation\big).
\end{equation*}

For a linear elastic material with Young's modulus $E$ and shear modulus $G$:
\begin{equation*}
\sigma = E\,\PointAxialStrain, \qquad \bm{\tau} = G\,\PointShearStrain.
\end{equation*}

The resultant reactions are related to the Ieşan parameters through a linear map:
\begin{equation}
\begin{pmatrix} 
N  \\ \ShearForce \\  T \\ \bm{M} 
\end{pmatrix}
= -\IesanMatrix \, \IesanVector
= -\begin{pmatrix}
EA & EA\CentroidLocation^\top & 0 & \mathbf{0}_{1\times2} \\
\mathbf{0}_{2\times1} & \mathbf{0}_{2\times2} & \mathbf{0}_{2\times1} & E\IntRGoRG \\
0 & \mathbf{0}_{1\times 2} & G\Jsv & \mathbf{0}_{1\times2} \\
EA \FrameBasis\times\CentroidLocation & \FrameBasis\times E\IntRoR & \mathbf{0}_{2\times 1} & \mathbf{0}_{2\times2} \\
\end{pmatrix}
\label{eq:resultant-map}
\end{equation}
where the \emph{resultant matrix} $\IesanMatrix \in \mathbb{R}^{6 \times 6}$ 
encodes the Saint-Venant constitutive relations  \eqref{eq:2-15} and \eqref{eq:2-18}.
$$
\IesanMatrix^{-1} =
\begin{pmatrix}
\displaystyle \frac{1}{EA} + \CentroidLocation\cdot (E\IntRGoRG)^{-1}\,\CentroidLocation
& \mathbf{0}_{1\times2}
& 0
& \CentroidLocation^\top (E\IntRGoRG)^{-1}\,\FrameBasis^\times
\\[6pt]
-(E\IntRGoRG)^{-1}\,\CentroidLocation
& \mathbf{0}_{2\times2}
& \mathbf{0}_{2\times1}
& -(E\IntRGoRG)^{-1}\,\FrameBasis^\times
\\[6pt]
0
& \mathbf{0}_{1\times2}
& \displaystyle \frac{1}{G\Jsv}
& \mathbf{0}_{1\times2}
\\[6pt]
\mathbf{0}_{2\times1}
& \bigl(E\IntRGoRG\bigr)^{-1}
& \mathbf{0}_{2\times1}
& \mathbf{0}_{2\times2}
\end{pmatrix}
$$

Thus the energetic trace is defined by the condition:
\[
\IesanMatrix \mathbf{X}^{-1} = \mathbf{X}^{-T}\underbrace{\int \mathbf{B}^T\mathbb{C}\mathbf{B}\dA}_{:=\bar{\mathbf{C}}} \mathbf{X}^{-1}
\quad\implies\quad 
\mathbf{X}^{-1} = \IesanMatrix^{-1} \bar{\mathbf{C}}^T
\]
with 
\[
\mathbb{C} := E\FrameBasis\otimes\FrameBasis + G\FrameBasis_{\alpha}\otimes\FrameBasis_{\alpha}
\]
Recall \eqref{eq:bishear-b-tensor}:
\begin{equation*}
\BishearTensorBB := \int\bm{W}_{(b)}^{\!\top}\,G\,\mathbf{B}_{(b)}\dA,
\qquad\text{and}\qquad 
\BishearTensorPB := \int\nabla\bm{\WarpShearIesan}\,G\,\mathbf{B}_{(b)}\dA
\end{equation*}

\paragraph{Work-conjugacy conditions.}

The virtual work of the section is:
\begin{equation}
\delta W = \int_\Section \delta\PointAxialStrain\,\sigma \, \mathrm{d}A 
+ \int_\Section \delta\PointShearStrain \cdot \bm{\tau} \, \mathrm{d}A.
\label{eq:virtual-work}
\end{equation}

We seek beam strains $(\epsilon, \bm{\gamma}_\perp, \varkappa, \bm{\kappa}_\perp)$ 
such that:
\begin{equation}
\delta W = \delta\epsilon\,N + \delta\bm{\gamma}_\perp \cdot \bm{F} 
+ \delta\varkappa\,T + \delta\bm{\kappa}_\perp \cdot \bm{M}.
\label{eq:conjugacy-target}
\end{equation}

We analyze each Ieşan mode separately.

\subparagraph{Extension.}

Consider $\delta a_\varepsilon \neq 0$ with all other parameters zero:
\begin{equation*}
\delta\PointAxialStrain = \delta a_\varepsilon, \qquad \delta\PointShearStrain = \mathbf{0}.
\end{equation*}
The virtual work is:
\begin{equation*}
\delta W = \delta a_\varepsilon \int_\Section \sigma \, \mathrm{d}A = \delta a_\varepsilon\,N.
\end{equation*}
Matching with \eqref{eq:conjugacy-target} requires:
\begin{equation*}
\boxed{\epsilon^{(\varepsilon)} = a_\varepsilon}
\end{equation*}

\subparagraph{Bending.}

Consider $\delta\bm{a}_\Omega \neq \mathbf{0}$ with all other parameters zero:
\begin{equation*}
\delta\PointAxialStrain = \bm{r} \cdot \delta\bm{a}_\Omega, \qquad \delta\PointShearStrain = \mathbf{0}.
\end{equation*}
The virtual work is:
\begin{equation*}
\delta W = \int_\Section (\bm{r} \cdot \delta\bm{a}_\Omega)\,\sigma \, \mathrm{d}A 
= \delta\bm{a}_\Omega \cdot \int_\Section \bm{r}\,\sigma \, \mathrm{d}A.
\end{equation*}

From the definition \eqref{eq:def-M}:
\begin{equation*}
\bm{M} = \FrameBasis \times \int_\Section \bm{r}\,\sigma \, \mathrm{d}A 
\quad\Longrightarrow\quad 
\int_\Section \bm{r}\,\sigma \, \mathrm{d}A = -\FrameBasis \times \bm{M} = \bm{M} \times \FrameBasis.
\end{equation*}
Thus:
\begin{equation*}
\delta W = \delta\bm{a}_\Omega \cdot (\bm{M} \times \FrameBasis) 
= (\delta\bm{a}_\Omega \times \bm{M}) \cdot \FrameBasis 
= -(\FrameBasis \times \delta\bm{a}_\Omega) \cdot \bm{M}.
\end{equation*}

For this to equal $\delta\bm{\kappa}_\perp \cdot \bm{M}$, we need:
\begin{equation*}
\boxed{\bm{\kappa}_\perp^{(a)} = -\FrameBasis \times \bm{a}_\Omega = -\FrameBasis^\times\,\bm{a}_\Omega}
\end{equation*}

\subparagraph{Torsion.}

Consider $\delta a_\vartheta \neq 0$ with all other parameters zero:
\begin{equation*}
\delta\PointAxialStrain = 0, \qquad 
\delta\PointShearStrain = \delta a_\vartheta\,(\FrameBasis \times \bm{r} + \nabla\WarpTwistIesan).
\end{equation*}
The virtual work is:
\begin{align*}
\delta W &= \int_\Section \delta a_\vartheta\,(\FrameBasis \times \bm{r} + \nabla\WarpTwistIesan) \cdot \bm{\tau} \, \mathrm{d}A \\
&= \delta a_\vartheta \int_\Section (\FrameBasis \times \bm{r}) \cdot \bm{\tau} \, \mathrm{d}A 
+ \delta a_\vartheta \int_\Section \nabla\WarpTwistIesan \cdot \bm{\tau} \, \mathrm{d}A.
\end{align*}

The first integral is $T$ by definition \eqref{eq:def-T}. For the second integral, 
we use the identity $\nabla \cdot (\WarpTwistIesan\,\bm{\tau}) = \nabla\WarpTwistIesan \cdot \bm{\tau} + \WarpTwistIesan\,\nabla \cdot \bm{\tau}$ 
and the divergence theorem:
\begin{equation*}
\int_\Section \nabla\WarpTwistIesan \cdot \bm{\tau} \, \mathrm{d}A 
= \oint_{\partial\Section} \WarpTwistIesan\,(\bm{\tau} \cdot \bm{\nu}) \, \mathrm{d}s 
- \int_\Section \WarpTwistIesan\,(\nabla \cdot \bm{\tau}) \, \mathrm{d}A = 0,
\end{equation*}
since $\bm{\tau} \cdot \bm{\nu} = 0$ on $\partial\Section$ (traction-free lateral surface) 
and $\nabla \cdot \bm{\tau} = 0$ (equilibrium in absence of body forces).

Thus $\delta W = \delta a_\vartheta\,T$, and matching requires:
\begin{equation}
\varkappa^{(\vartheta)} = a_\vartheta
\label{eq:conj-twist-torsion}
\end{equation}

\subparagraph{Flexure.}

This is the most involved case. Consider $\delta\bm{b}_{\Section} \neq \mathbf{0}$ with 
all other parameters zero:
\begin{align}
\delta\PointAxialStrain &= x\,(\bm{r} - \CentroidLocation) \cdot \delta\bm{b}_{\Section}, 
\label{eq:var-axial-flexure} \\
\delta\PointShearStrain &= \delta c_\vartheta\,(\FrameBasis \times \bm{r} + \nabla\WarpTwistIesan) 
+ \bm{W}_{(b)}\,\delta\bm{b}_{\Section} + (\nabla\bm{\WarpShearIesan})^\top\,\delta\bm{b}_{\Section},
\label{eq:var-shear-flexure}
\end{align}
where $\delta c_\vartheta = \CenterShearLocation \cdot \delta\bm{b}_{\Section}$.

\emph{Axial contribution.} From \eqref{eq:var-axial-flexure}:
\begin{align*}
\delta W_\sigma &= \int_\Section x\,(\bm{r} - \CentroidLocation) \cdot \delta\bm{b}_{\Section}\,\sigma \, \mathrm{d}A \\
&= x\,\delta\bm{b}_{\Section} \cdot \int_\Section (\bm{r} - \CentroidLocation)\,\sigma \, \mathrm{d}A \\
&= x\,\delta\bm{b}_{\Section} \cdot \left( \int_\Section \bm{r}\,\sigma \, \mathrm{d}A 
- \CentroidLocation \int_\Section \sigma \, \mathrm{d}A \right) \\
&= x\,\delta\bm{b}_{\Section} \cdot \left( \bm{M} \times \FrameBasis - \CentroidLocation\,N \right).
\end{align*}

Rearranging:
\begin{equation}
\delta W_\sigma = -x\,(\FrameBasis \times \delta\bm{b}_{\Section}) \cdot \bm{M} 
- x\,(\CentroidLocation \cdot \delta\bm{b}_{\Section})\,N.
\label{eq:axial-work-flexure}
\end{equation}

\emph{Shear contribution.} From \eqref{eq:var-shear-flexure}:
\begin{align}
\delta W_\tau &= \int_\Section \delta\PointShearStrain \cdot \bm{\tau} \, \mathrm{d}A \notag \\
&= \delta c_\vartheta \int_\Section (\FrameBasis \times \bm{r} + \nabla\WarpTwistIesan) \cdot \bm{\tau} \, \mathrm{d}A 
+ \delta\bm{b}_{\Section} \cdot \int_\Section \bm{W}_{(b)}^\top\,\bm{\tau} \, \mathrm{d}A \notag \\
&\quad + \delta\bm{b}_{\Section} \cdot \int_\Section \nabla\bm{\WarpShearIesan}\,\bm{\tau} \, \mathrm{d}A.
\label{eq:shear-work-flexure-1}
\end{align}

The first integral: By the same argument as for torsion (using $\bm{\tau} \cdot \bm{\nu} = 0$ 
and $\nabla \cdot \bm{\tau} = 0$), the $\nabla\WarpTwistIesan$ term vanishes:
\begin{equation*}
\int_\Section (\FrameBasis \times \bm{r} + \nabla\WarpTwistIesan) \cdot \bm{\tau} \, \mathrm{d}A = T.
\end{equation*}

The third integral: Using integration by parts on each column of $\bm{\WarpShearIesan}$:
\begin{equation*}
\int_\Section \nabla\WarpShearIesan_\alpha\,\bm{\tau} \, \mathrm{d}A 
= \oint_{\partial\Section} \WarpShearIesan_\alpha\,(\bm{\tau} \cdot \bm{\nu}) \, \mathrm{d}s 
- \int_\Section \WarpShearIesan_\alpha\,(\nabla \cdot \bm{\tau}) \, \mathrm{d}A = \mathbf{0}.
\end{equation*}
Thus $\int_\Section \nabla\bm{\WarpShearIesan}\,\bm{\tau} \, \mathrm{d}A = \mathbf{0}$.

Substituting back:
\begin{equation}
\delta W_\tau = (\CenterShearLocation \cdot \delta\bm{b}_{\Section})\,T 
+ \delta\bm{b}_{\Section} \cdot \int_\Section \bm{W}_{(b)}^\top\,\bm{\tau} \, \mathrm{d}A.
\label{eq:shear-work-flexure-2}
\end{equation}

\emph{Total virtual work from flexure.} Combining \eqref{eq:axial-work-flexure} 
and \eqref{eq:shear-work-flexure-2}:
\begin{align}
\delta W^{(b)} &= -x\,(\CentroidLocation \cdot \delta\bm{b}_{\Section})\,N 
- x\,(\FrameBasis \times \delta\bm{b}_{\Section}) \cdot \bm{M} \notag \\
&\quad + (\CenterShearLocation \cdot \delta\bm{b}_{\Section})\,T 
+ \delta\bm{b}_{\Section} \cdot \int_\Section \bm{W}_{(b)}^\top\,\bm{\tau} \, \mathrm{d}A.
\label{eq:total-work-flexure}
\end{align}

\emph{Matching with conjugacy condition.} We require:
\begin{equation}
\delta W^{(b)} = \delta\epsilon^{(b)}\,N + \delta\bm{\gamma}_\perp^{(b)} \cdot \bm{F} 
+ \delta\varkappa^{(b)}\,T + \delta\bm{\kappa}_\perp^{(b)} \cdot \bm{M}.
\label{eq:match-flexure}
\end{equation}

Comparing term by term:
\begin{itemize}
\item Coefficient of $N$: $\delta\epsilon^{(b)} = -x\,(\CentroidLocation \cdot \delta\bm{b}_{\Section})$
\item Coefficient of $\bm{M}$: $\delta\bm{\kappa}_\perp^{(b)} = -x\,(\FrameBasis \times \delta\bm{b}_{\Section})$
\item Coefficient of $T$: $\delta\varkappa^{(b)} = \CenterShearLocation \cdot \delta\bm{b}_{\Section}$
\item Coefficient of $\bm{F}$: This requires careful treatment of the remaining term.
\end{itemize}

The remaining term is $\delta\bm{b}_{\Section} \cdot \int_\Section \bm{W}_{(b)}^\top\,\bm{\tau} \, \mathrm{d}A$. 
For this to equal $\delta\bm{\gamma}_\perp^{(b)} \cdot \bm{F}$, we need to express 
it in terms of the shear force $\bm{F} = \int_\Section \bm{\tau} \, \mathrm{d}A$.

Define the \emph{Poisson-shear tensor}:
\begin{equation}
\mathbf{P} := \int_\Section \bm{W}_{(b)}^\top\,\bm{\tau} \, \mathrm{d}A 
\quad \in \mathbb{R}^2,
\label{eq:def-P}
\end{equation}
which depends on the stress state $\bm{\tau}$. Under flexure loading, 
$\bm{\tau} = G\,\PointShearStrain^{(b)}$ where:
\begin{equation*}
\PointShearStrain^{(b)} = c_\vartheta\,(\FrameBasis \times \bm{r} + \nabla\WarpTwistIesan) 
+ \bm{W}_{(b)}\,\bm{b}_{\Section} + (\nabla\bm{\WarpShearIesan})^\top\,\bm{b}_{\Section}.
\end{equation*}

The shear force under flexure is:
\begin{align}
\bm{F} &= G \int_\Section \PointShearStrain^{(b)} \, \mathrm{d}A \notag \\
&= G\,c_\vartheta \int_\Section (\FrameBasis \times \bm{r} + \nabla\WarpTwistIesan) \, \mathrm{d}A 
+ G \int_\Section \bm{W}_{(b)} \, \mathrm{d}A\,\bm{b}_{\Section} 
+ G \int_\Section (\nabla\bm{\WarpShearIesan})^\top \, \mathrm{d}A\,\bm{b}_{\Section}.
\label{eq:F-flexure}
\end{align}

Using $\int_\Section (\FrameBasis \times \bm{r}) \, \mathrm{d}A = A\,(\FrameBasis \times \CentroidLocation)$ 
and noting that $\int_\Section \nabla\WarpTwistIesan \, \mathrm{d}A = \oint_{\partial\Section} \WarpTwistIesan\,\bm{\nu} \, \mathrm{d}s$ 
(which need not vanish), define:
\begin{align}
\bm{q}_\varphi &:= \int_\Section (\FrameBasis \times \bm{r} + \nabla\WarpTwistIesan) \, \mathrm{d}A, 
\label{eq:def-q-phi} \\
\mathbf{Q}_W &:= \int_\Section \bm{W}_{(b)} \, \mathrm{d}A = A\,\SectionAverage{\bm{W}_{(b)}}, 
\label{eq:def-Q-W} \\
\mathbf{Q}_\psi &:= \int_\Section (\nabla\bm{\WarpShearIesan})^\top \, \mathrm{d}A.
\label{eq:def-Q-psi}
\end{align}

Then:
\begin{equation}
\bm{F} = G\,c_\vartheta\,\bm{q}_\varphi + G\,(\mathbf{Q}_W + \mathbf{Q}_\psi)\,\bm{b}_{\Section} 
= G\,(\CenterShearLocation \cdot \bm{b}_{\Section})\,\bm{q}_\varphi + G\,\mathbf{Q}\,\bm{b}_{\Section},
\label{eq:F-flexure-compact}
\end{equation}
where $\mathbf{Q} := \mathbf{Q}_W + \mathbf{Q}_\psi \in \mathbb{R}^{3 \times 2}$.

\emph{The shear conjugacy problem.} We need $\delta\bm{\gamma}_\perp^{(b)} \cdot \bm{F}$ 
to equal $\delta\bm{b}_{\Section} \cdot \mathbf{P}$. 
This is a compatibility condition 
between the shear strain definition and the Poisson-shear coupling.

In general, $\mathbf{P}$ and $\bm{F}$ are not simply related by a constant tensor 
acting on $\delta\bm{b}_{\Section}$. However, we can define the energy-conjugate shear 
strain implicitly.

Write:
\begin{equation}
\delta\bm{b}_{\Section} \cdot \mathbf{P} = \delta\bm{b}_{\Section} \cdot \mathbf{P}(\bm{b}_{\Section}),
\end{equation}
where $\mathbf{P}(\bm{b}_{\Section})$ is linear in $\bm{b}_{\Section}$ through the stress $\bm{\tau}$.

For the constitutive relation to be symmetric, we require:
\begin{equation}
\bm{\gamma}_\perp^{(b)} = \mathbf{K}^{-1}\,\bm{F},
\label{eq:romano-shear-def}
\end{equation}
where $\mathbf{K}$ is the shear stiffness tensor defined by $\bm{F} = \mathbf{K}\,\bm{\gamma}_\perp$.

From \eqref{eq:F-flexure-compact}, the shear force is linear in $\bm{b}_{\Section}$:
\begin{equation}
\bm{F} = \mathbf{R}\,\bm{b}_{\Section}, \qquad 
\mathbf{R} := G\,(\bm{q}_\varphi \otimes \CenterShearLocation + \mathbf{Q}).
\label{eq:F-R-relation}
\end{equation}

For energy conjugacy with symmetric stiffness, Romano defines:
\begin{equation}
\bm{\gamma}_\perp^{(b)} := \mathbf{K}^{-1}\,\mathbf{R}\,\bm{b}_{\Section},
\label{eq:ec-shear-final}
\end{equation}
where the shear stiffness $\mathbf{K}$ is determined by the requirement that 
the strain energy $\frac{1}{2}\bm{\gamma}_\perp \cdot \bm{F}$ equals the actual 
shear strain energy $\frac{1}{2}\int_\Section \PointShearStrain \cdot \bm{\tau} \, \mathrm{d}A$.

This gives:
\begin{equation}
\mathbf{K} = \mathbf{R}\,\mathbf{H}^{-1}\,\mathbf{R}^\top,
\label{eq:K-formula}
\end{equation}
where $\mathbf{H}$ is the shear energy tensor:
\begin{equation}
\mathbf{H} := G \int_\Section \PointShearStrain^{(b)} \otimes \PointShearStrain^{(b)} \, \mathrm{d}A 
\bigg|_{\bm{b}_{\Section} = \text{unit}}.
\label{eq:def-H}
\end{equation}

Substituting \eqref{eq:K-formula} into \eqref{eq:ec-shear-final}:
\begin{equation}
\boxed{\bm{\gamma}_\perp^{(b)} = \mathbf{H}^{-1}\,\mathbf{R}^\top\,(\mathbf{R}\,\mathbf{H}^{-1}\,\mathbf{R}^\top)^{-1}\,\mathbf{R}\,\bm{b}_{\Section} 
= \mathbf{H}^{-1}\,\mathbf{R}^\top\,\bm{\lambda}}
\label{eq:ec-shear-explicit}
\end{equation}
where $\bm{\lambda}$ is a Lagrange multiplier enforcing the force-strain relation.

For the simpler case where we directly define shear strain from the force:
\begin{equation}
\boxed{\bm{\gamma}_\perp^{(b)} = \mathbf{K}^{-1}\,\bm{F} = \mathbf{K}^{-1}\,\mathbf{R}\,\bm{b}_{\Section}}
\end{equation}

\paragraph{Summary of energy-conjugate strains.}

Collecting results from all modes:
\begin{align}
\epsilon &= a_\varepsilon - x\,(\CentroidLocation \cdot \bm{b}_{\Section}), 
\label{eq:ec-axial-final} \\
\bm{\gamma}_\perp &= \mathbf{K}^{-1}\,\mathbf{R}\,\bm{b}_{\Section}, 
\label{eq:ec-shear-final2} \\
\varkappa &= a_\vartheta + \CenterShearLocation \cdot \bm{b}_{\Section}, 
\label{eq:ec-twist-final} \\
\bm{\kappa}_\perp &= -\FrameBasis^\times\,\bm{a}_\Omega - x\,(\FrameBasis \times \bm{b}_{\Section}).
\label{eq:ec-curv-final}
\end{align}

\paragraph{Assembly of trace matrices.}

The beam strains $\bm{e}(x) = (\epsilon, \bm{\gamma}_\perp, \varkappa, \bm{\kappa}_\perp)^\top$ 
satisfy:
\begin{equation}
\bm{e}(x) = \TraceDeDpRoot\,\bm{p} + x\,\TraceDeDpOne\,\bm{p},
\end{equation}
with $\bm{p} = (a_\varepsilon, \bm{a}_\Omega, a_\vartheta, \bm{b}_{\Section})^\top$ and:

\begin{equation}
\TraceDeDpRoot = \begin{pmatrix}
1 & \mathbf{0}^\top & 0 & \mathbf{0}^\top \\[4pt]
\mathbf{0} & \mathbf{0} & \mathbf{0} & \mathbf{K}^{-1}\mathbf{R} \\[4pt]
0 & \mathbf{0}^\top & 1 & \CenterShearLocation^\top \\[4pt]
\mathbf{0} & -\FrameBasis^\times & \mathbf{0} & \mathbf{0}
\end{pmatrix}
\label{eq:T0-romano-final}
\end{equation}

\begin{equation}
\TraceDeDpOne = \begin{pmatrix}
0 & \mathbf{0}^\top & 0 & -\CentroidLocation^\top \\[4pt]
\mathbf{0} & \mathbf{0} & \mathbf{0} & \mathbf{0} \\[4pt]
0 & \mathbf{0}^\top & 0 & \mathbf{0}^\top \\[4pt]
\mathbf{0} & \mathbf{0} & \mathbf{0} & -\FrameBasis^\times
\end{pmatrix}
\label{eq:T1-romano-final}
\end{equation}

\paragraph{Warping index.}

The stacked matrix is:
\begin{equation}
\mathbf{T}_{\rm stack} = \begin{pmatrix} \TraceDeDpRoot \\ \TraceDeDpOne \end{pmatrix} 
\in \mathbb{R}^{12 \times 6}.
\end{equation}

Column analysis:
\begin{itemize}
\item $a_\varepsilon$: Column $(1, \mathbf{0}, 0, \mathbf{0}; 0, \mathbf{0}, 0, \mathbf{0})^\top$ — rank 1.
\item $\bm{a}_\Omega$: Columns with $-\FrameBasis^\times$ in curvature block — rank 2.
\item $a_\vartheta$: Column $(0, \mathbf{0}, 1, \mathbf{0}; 0, \mathbf{0}, 0, \mathbf{0})^\top$ — rank 1.
\item $\bm{b}_{\Section}$: In $\TraceDeDpRoot$: $\mathbf{K}^{-1}\mathbf{R}$ (shear) and $\CenterShearLocation^\top$ (twist). 
In $\TraceDeDpOne$: $-\CentroidLocation^\top$ (axial) and $-\FrameBasis^\times$ (curvature).
\end{itemize}

The $\bm{b}_{\Section}$ columns contribute rank 2 through $\TraceDeDpOne$ (via $-\FrameBasis^\times$), 
independent of the contributions from $\bm{a}_\Omega$ (which appear in different rows 
of the stacked matrix).

\begin{proposition}[Warping index of energy-conjugate trace]
\label{prop:warp-index-romano-final}
For sections with non-degenerate geometry, the energy-conjugate trace has 
warping index:
\begin{equation}
w = 6 - \mathrm{rank}(\mathbf{T}_{\rm stack}) = 0.
\end{equation}
\end{proposition}

\begin{proof}
The six columns of $\mathbf{T}_{\rm stack}$ are:
\begin{enumerate}
\item $a_\varepsilon$: Nonzero in row 1 of $\TraceDeDpRoot$.
\item $a_{\Omega,1}$: Nonzero in rows 5--6 of $\TraceDeDpRoot$ (via $-\FrameBasis^\times$).
\item $a_{\Omega,2}$: Nonzero in rows 5--6 of $\TraceDeDpRoot$ (via $-\FrameBasis^\times$).
\item $a_\vartheta$: Nonzero in row 3 of $\TraceDeDpRoot$.
\item $b_{\Omega,1}$: Nonzero in row 1 of $\TraceDeDpOne$ (via $-\CentroidLocation^\top$) 
and rows 11--12 of $\TraceDeDpOne$ (via $-\FrameBasis^\times$).
\item $b_{\Omega,2}$: Same structure as $b_{\Omega,1}$.
\end{enumerate}
These columns are linearly independent for generic $\CentroidLocation \neq \mathbf{0}$ 
(non-centroidal origin), giving $\mathrm{rank}(\mathbf{T}_{\rm stack}) = 6$.

For centroidal coordinates ($\CentroidLocation = \mathbf{0}$), the $\bm{b}_{\Section}$ columns 
still contribute through $-\FrameBasis^\times$ in $\TraceDeDpOne$ and $\mathbf{K}^{-1}\mathbf{R}$ 
in $\TraceDeDpRoot$, maintaining full rank.
\end{proof}

\begin{remark}[Constitutive symmetry]
By construction, the energy-conjugate trace yields a symmetric section 
stiffness matrix:
\begin{equation}
\bm{s} = \mathbf{C}\,\bm{e}, \qquad \mathbf{C} = \mathbf{C}^\top,
\end{equation}
where $\bm{s} = (N, \bm{F}, T, \bm{M})^\top$. This follows from the variational 
derivation: the strains are defined as derivatives of the complementary 
energy with respect to the resultants.
\end{remark}

\begin{remark}[Comparison with section-average trace]
The two traces yield identical results for:
\begin{itemize}
\item Axial strain from extension: $\epsilon^{(\varepsilon)} = a_\varepsilon$
\item Curvature from bending: $\bm{\kappa}_\perp^{(a)} = -\FrameBasis^\times\,\bm{a}_\Omega$
\item Twist from torsion: $\varkappa^{(\vartheta)} = a_\vartheta$
\item Curvature from flexure: $\bm{\kappa}_\perp^{(b)} = -x\,(\FrameBasis \times \bm{b}_{\Section})$
\end{itemize}
They differ in:
\begin{itemize}
\item Shear strain from flexure: section-average gives $\bm{\Gamma}_0\,\bm{b}_{\Section}$; 
energy-conjugate gives $\mathbf{K}^{-1}\mathbf{R}\,\bm{b}_{\Section}$.
\item Twist from flexure: section-average includes additional coupling terms; 
energy-conjugate gives $\CenterShearLocation \cdot \bm{b}_{\Section}$.
\end{itemize}
The difference is a gauge transformation that preserves the warping index 
but changes the constitutive coupling structure.
\end{remark}

\begin{remark}[Plane-section property]
The energy-conjugate trace does not satisfy the plane-section property. 
For a section-rigid displacement $\hat{\bm{u}} = \bm{a} + \bm{\omega} \times \bm{r}$, 
the extracted beam strains $(\bm{\gamma}, \bm{\kappa})$ are not generally 
$(\bm{a}' + \FrameBasis \times \bm{\omega}, \bm{\omega}')$ evaluated at a consistent 
reference point. 
\end{remark}
% \clearpage
% 
\textbf{Setup.} The local parameters are
\(\hat{\bm{e}} = (\epsilon_0, \bm{\beta}, \tau, \bm{\epsilon}_1)\) with
sizes \((1, 2, 1, 2)\). The 3D strain is:
\[\bm{\varepsilon} = \varepsilon\,\FrameBasis \otimes \FrameBasis + \FrameBasis \stackrel{s}{\otimes} \bm{\varepsilon}_\gamma - \nu\varepsilon\,\mathbf{1}_\Omega\]
where:
\[\varepsilon = \epsilon_0 + \bm{r} \cdot \bm{\epsilon}_1, \qquad \bm{\varepsilon}_\gamma = \tau\,\bm{\tau}_\varphi + \bm{\tau}_\psi\,\bm{\beta}\]
with:
\[\bm{\tau}_\varphi := \FrameBasis \times \bm{r} + \nabla\WarpTwistIesan, \qquad \bm{\tau}_\psi := \nabla\bm{\WarpShearIesan} + \bm{\Pi}_{(b)}\]

\textbf{Computing
\(\bar{\mathbf{C}} = \int_\Section \mathbf{B}^T \mathbb{C} \mathbf{B} \dA\).}

The strain energy density splits into axial and shear parts:
\[\tfrac{1}{2}\bm{\varepsilon}^T \mathbb{C} \bm{\varepsilon} = \tfrac{1}{2}E\varepsilon^2 + \tfrac{1}{2}G|\bm{\varepsilon}_\gamma|^2\]

In the partition \((\epsilon_0, \bm{\beta}, \tau, \bm{\epsilon}_1)\):

\[\bar{\mathbf{C}} = \begin{pmatrix}
EA & \mathbf{0}^\top & 0 & EA\CentroidLocation^\top \\[4pt]
\mathbf{0} & G\mathbf{H}_{\psi\psi} & G\mathbf{h}_{\varphi\psi} & \mathbf{0} \\[4pt]
0 & G\mathbf{h}_{\varphi\psi}^\top & G\Jsv & \mathbf{0}^\top \\[4pt]
EA\CentroidLocation & \mathbf{0} & \mathbf{0} & E\IntRoR
\end{pmatrix}\]

where:
\[\Jsv := \int_\Section |\bm{\tau}_\varphi|^2 \dA, \qquad \mathbf{h}_{\varphi\psi} := \int_\Section \bm{\tau}_\psi^\top \bm{\tau}_\varphi \dA, \qquad \mathbf{H}_{\psi\psi} := \int_\Section \bm{\tau}_\psi^\top \bm{\tau}_\psi \dA\]

\textbf{Evaluating \(\mathbf{h}_{\varphi\psi}\).}

We have:
\[\mathbf{h}_{\varphi\psi} = \int_\Section (\nabla\bm{\WarpShearIesan} + \bm{\Pi}_{(b)})^\top (\FrameBasis \times \bm{r} + \nabla\WarpTwistIesan) \dA\]

Breaking into four terms:

\emph{Term 1:}
\(\int_\Section (\nabla\bm{\WarpShearIesan})^\top (\FrameBasis \times \bm{r}) \dA\)

Using divergence theorem with
\(\nabla \cdot (\FrameBasis \times \bm{r}) = 0\):
\[\int_\Section (\FrameBasis \times \bm{r}) \cdot \nabla\WarpShearIesan_\alpha \dA = \oint_{\partial\Section} \WarpShearIesan_\alpha (\FrameBasis \times \bm{r}) \cdot \bm{\nu} \, ds\]

\emph{Term 2:}
\(\int_\Section (\nabla\bm{\WarpShearIesan})^\top \nabla\WarpTwistIesan \dA\)

Using Green's identity with \(\Delta\WarpTwistIesan = 0\) and
\(\partial_\nu\WarpTwistIesan = -(\FrameBasis \times \bm{r}) \cdot \bm{\nu}\):
\[\int_\Section \nabla\WarpTwistIesan \cdot \nabla\WarpShearIesan_\alpha \dA = \oint_{\partial\Section} \WarpShearIesan_\alpha \partial_\nu\WarpTwistIesan \, ds = -\oint_{\partial\Section} \WarpShearIesan_\alpha (\FrameBasis \times \bm{r}) \cdot \bm{\nu} \, ds\]

Terms 1 and 2 cancel:
\[\int_\Section (\nabla\bm{\WarpShearIesan})^\top (\FrameBasis \times \bm{r} + \nabla\WarpTwistIesan) \dA = \mathbf{0}\]

\emph{Term 3:}
\(\int_\Section \bm{\Pi}_{(b)}^\top (\FrameBasis \times \bm{r}) \dA\)

With
\(\bm{\Pi}_{(b)} = -\nu(\text{dev}(\bm{r} \otimes \bm{r}) - \bm{r} \otimes \CentroidLocation)\),
acting on column \(\alpha\):
\[(\bm{\Pi}_{(b)})_\alpha = -\nu\big((r_\alpha - \bar{\rho}_\alpha)\bm{r} - \tfrac{1}{2}|\bm{r}|^2 \bm{e}_\alpha\big)\]

Then:
\[(\FrameBasis \times \bm{r}) \cdot (\bm{\Pi}_{(b)})_\alpha = -\nu(r_\alpha - \bar{\rho}_\alpha)\underbrace{(\FrameBasis \times \bm{r}) \cdot \bm{r}}_{=0} + \tfrac{\nu}{2}|\bm{r}|^2 (\FrameBasis \times \bm{r})_\alpha\]

So:
\[\int_\Section \bm{\Pi}_{(b)}^\top (\FrameBasis \times \bm{r}) \dA = \tfrac{\nu}{2} \int_\Section |\bm{r}|^2 (\FrameBasis \times \bm{r}) \dA\]

\emph{Term 4:}
\(\int_\Section \bm{\Pi}_{(b)}^\top \nabla\WarpTwistIesan \dA\)

\[\int_\Section (\bm{\Pi}_{(b)})_\alpha \cdot \nabla\WarpTwistIesan \dA = -\nu \int_\Section (r_\alpha - \bar{\rho}_\alpha)(\bm{r} \cdot \nabla\WarpTwistIesan) \dA + \tfrac{\nu}{2} \int_\Section |\bm{r}|^2 \partial_\alpha\WarpTwistIesan \dA\]

Combining terms 3 and 4:
\[\mathbf{h}_{\varphi\psi} = \tfrac{\nu}{2} \int_\Section |\bm{r}|^2 (\FrameBasis \times \bm{r} + \nabla\WarpTwistIesan) \dA - \nu \int_\Section (\bm{r} - \CentroidLocation)(\bm{r} \cdot \nabla\WarpTwistIesan) \dA\]

Define:
\[\bm{\mu}_\varphi := \int_\Section |\bm{r}|^2 \bm{\tau}_\varphi \dA, \qquad \bm{\lambda}_\varphi := \int_\Section (\bm{r} \cdot \nabla\WarpTwistIesan)(\bm{r} - \CentroidLocation) \dA\]

Then:
\[\mathbf{h}_{\varphi\psi} = \tfrac{\nu}{2}\bm{\mu}_\varphi - \nu\bm{\lambda}_\varphi\]

\textbf{Relation to shear center.} The shear center is defined by:
\[\Jsv\,\CenterShearLocation = -\int_\Section \bm{r} \times \bm{\tau}_\psi \cdot \FrameBasis \dA = \int_\Section (\FrameBasis \times \bm{r})^\top \bm{\tau}_\psi \dA\]

This is \emph{not} the same as \(\mathbf{h}_{\varphi\psi}\) because the
latter includes the \(\nabla\WarpTwistIesan\) term.

\subsubsection{The energy-conjugacy
condition.}\label{the-energy-conjugacy-condition.}

For beam strains \(\bm{e} = \mathbf{X}\hat{\bm{e}}\) to be
work-conjugate to resultants
\(\bm{s} = \mathbf{R}\mathbf{L}^{-1}\hat{\bm{e}}\):
\[\bm{s}^\top \delta\bm{e} = \int_\Section \bm{\sigma}^\top \delta\bm{\varepsilon} \dA = \hat{\bm{e}}^\top \bar{\mathbf{C}} \delta\hat{\bm{e}}\]

This requires:
\[\mathbf{X}^\top \mathbf{R} \mathbf{L}^{-1} = \bar{\mathbf{C}}\]

Hence:
\[\mathbf{X}^{-1} = \bar{\mathbf{C}}^{-1} \mathbf{L}^{-\top} \mathbf{R}^\top\]

\textbf{Computing \(\mathbf{L}^{-\top} \mathbf{R}^\top\).}

Recall (with partition
\(\hat{\bm{e}} = (\epsilon_0, \bm{\beta}, \tau, \bm{\epsilon}_1)\) and
\(\bm{p} = (a_\varepsilon, \bm{a}_\Omega, a_\vartheta, \bm{b}_{\Section})\)):

\[\mathbf{L}_0 = \begin{pmatrix}
1 & \mathbf{0}^\top & 0 & \mathbf{0}^\top \\
\mathbf{0} & \mathbf{0} & \mathbf{0} & \oneS \\
0 & \mathbf{0}^\top & 1 & \CenterShearLocation^\top \\
\mathbf{0} & \oneS & \mathbf{0} & \mathbf{0}
\end{pmatrix}
\qquad\text{so}\qquad 
\mathbf{L}_0^{-\top} = \begin{pmatrix}
1 & \mathbf{0}^\top & 0 & \mathbf{0}^\top \\
\mathbf{0} & \mathbf{0} & -\CenterShearLocation & \oneS \\
0 & \mathbf{0}^\top & 1 & \mathbf{0}^\top \\
\mathbf{0} & \oneS & \mathbf{0} & \mathbf{0}
\end{pmatrix}\]

And with partition \(\bm{s} = (N, \bm{F}, T, \bm{M})\) (but noting
\(\bm{M}\) appears as \(\FrameBasis \times \bm{M}\) in the formulas):

\[\mathbf{R}^\top = \begin{pmatrix}
-EA & \mathbf{0}^\top & 0 & -EA(\FrameBasis \times \CentroidLocation)^\top \\
-EA\CentroidLocation & \mathbf{0} & \mathbf{0} & -(\FrameBasis \times E\IntRoR)^\top \\
0 & \mathbf{0}^\top & -G\Jsv & \mathbf{0}^\top \\
\mathbf{0} & -E\IntRGoRG & \mathbf{0} & \mathbf{0}
\end{pmatrix}\]


\paragraph{Energy matrix \(\bar{\mathbf{C}}\).}

In partition
\(\hat{\bm{e}} = (\epsilon_0, \bm{\beta}, \tau, \bm{\epsilon}_1)\):
\[\bar{\mathbf{C}} = \begin{pmatrix}
EA & \mathbf{0}^\top & 0 & EA\CentroidLocation^\top \\
\mathbf{0} & G\mathbf{H}_{\psi\psi} & G\mathbf{h}_{\varphi\psi} & \mathbf{0} \\
0 & G\mathbf{h}_{\varphi\psi}^\top & G\Jsv & \mathbf{0}^\top \\
EA\CentroidLocation & \mathbf{0} & \mathbf{0} & E\IntRoR
\end{pmatrix}
\]

\begin{Details}
Verifying the \((3,3)\) entry (torsional stiffness):

\[
\int_\Section |\bm{\tau}_\varphi|^2 \dA = \int_\Section |\FrameBasis \times \bm{r} + \nabla\WarpTwistIesan|^2 \dA
\]
\[= \int_\Section r^2 \dA + 2\int_\Section (\FrameBasis \times \bm{r}) \cdot \nabla\WarpTwistIesan \dA + \int_\Section |\nabla\WarpTwistIesan|^2 \dA
\]

For the last term, using Green's first identity with
\(\Delta\WarpTwistIesan = 0\):
\[\int_\Section |\nabla\WarpTwistIesan|^2 \dA = \oint_{\partial\Section} \WarpTwistIesan \,\partial_\nu\WarpTwistIesan \, ds = -\oint_{\partial\Section} \WarpTwistIesan (\FrameBasis \times \bm{r}) \cdot \bm{\nu} \, ds\]

By divergence theorem (using
\(\nabla \cdot (\FrameBasis \times \bm{r}) = 0\)):
\[\oint_{\partial\Section} \WarpTwistIesan (\FrameBasis \times \bm{r}) \cdot \bm{\nu} \, ds = \int_\Section \nabla\WarpTwistIesan \cdot (\FrameBasis \times \bm{r}) \dA\]

Therefore:
\[\int_\Section |\nabla\WarpTwistIesan|^2 \dA = -\int_\Section (\FrameBasis \times \bm{r}) \cdot \nabla\WarpTwistIesan \dA\]

So:
\[\int_\Section |\bm{\tau}_\varphi|^2 \dA = \int_\Section r^2 \dA + \int_\Section (\FrameBasis \times \bm{r}) \cdot \nabla\WarpTwistIesan \dA = \Jsv\]

\textbf{The (2,3) entry (shear-torsion coupling):}

\[\mathbf{h}_{\varphi\psi} = \int_\Section \bm{\tau}_\psi^\top \bm{\tau}_\varphi \dA = \int_\Section (\bm{\Pi}_{(b)}^\top + \nabla\bm{\WarpShearIesan})(\FrameBasis \times \bm{r} + \nabla\WarpTwistIesan) \dA\]

Breaking into four terms:

\emph{Term A:}
\(\int_\Section (\nabla\bm{\WarpShearIesan})(\FrameBasis \times \bm{r}) \dA\)

\emph{Term B:}
\(\int_\Section (\nabla\bm{\WarpShearIesan})\nabla\WarpTwistIesan \dA\)

Using Green's identity with \(\Delta\WarpTwistIesan = 0\) and
\(\partial_\nu\WarpTwistIesan = -(\FrameBasis \times \bm{r}) \cdot \bm{\nu}\):
\[[\text{Term B}]_\alpha = \oint_{\partial\Section} \WarpShearIesan_\alpha \,\partial_\nu\WarpTwistIesan \, ds = -\oint_{\partial\Section} \WarpShearIesan_\alpha (\FrameBasis \times \bm{r}) \cdot \bm{\nu} \, ds\]

By divergence theorem:
\[[\text{Term A}]_\alpha = \oint_{\partial\Section} \WarpShearIesan_\alpha (\FrameBasis \times \bm{r}) \cdot \bm{\nu} \, ds\]

So Terms A and B cancel:
\(\int_\Section (\nabla\bm{\WarpShearIesan})\bm{\tau}_\varphi \dA = \mathbf{0}\).
✓

\emph{Term C:}
\(\int_\Section \bm{\Pi}_{(b)}^\top (\FrameBasis \times \bm{r}) \dA\)

\emph{Term D:}
\(\int_\Section \bm{\Pi}_{(b)}^\top \nabla\WarpTwistIesan \dA\)

For the \(\alpha\)-th component, with
\((\bm{\Pi}_{(b)})_\alpha = -\nu((r_\alpha - \bar{\rho}_\alpha)\bm{r} - \frac{1}{2}r^2\bm{e}_\alpha)\):

\textbf{Term C:}
\[(\bm{\Pi}_{(b)})_\alpha \cdot (\FrameBasis \times \bm{r}) = -\nu(r_\alpha - \bar{\rho}_\alpha)\underbrace{\bm{r} \cdot (\FrameBasis \times \bm{r})}_{=0} + \frac{\nu}{2}r^2(\FrameBasis \times \bm{r})_\alpha\]

So:
\[\text{Term C} = \frac{\nu}{2}\int_\Section r^2 (\FrameBasis \times \bm{r}) \dA\]

\textbf{Term D:}
\[(\bm{\Pi}_{(b)})_\alpha \cdot \nabla\WarpTwistIesan = -\nu(r_\alpha - \bar{\rho}_\alpha)(\bm{r} \cdot \nabla\WarpTwistIesan) + \frac{\nu}{2}r^2 \partial_\alpha\WarpTwistIesan\]

So:
\[\text{Term D} = \frac{\nu}{2}\int_\Section r^2 \nabla\WarpTwistIesan \dA - \nu\int_\Section (\bm{r} - \CentroidLocation)(\bm{r} \cdot \nabla\WarpTwistIesan) \dA\]

Therefore:
\[\mathbf{h}_{\varphi\psi} = \frac{\nu}{2}\int_\Section r^2 \bm{\tau}_\varphi \dA - \nu\int_\Section (\bm{r} - \CentroidLocation)(\bm{r} \cdot \nabla\WarpTwistIesan) \dA\]

\textbf{Relation to shear center.} The shear center satisfies:
\[\Jsv\,\CenterShearLocation = \int_\Section \bm{\tau}_\psi^\top (\bm{r} \times \FrameBasis) \dA = -\int_\Section \bm{\tau}_\psi^\top (\FrameBasis \times \bm{r}) \dA\]

This involves only \((\FrameBasis \times \bm{r})\), not the full
\(\bm{\tau}_\varphi = \FrameBasis \times \bm{r} + \nabla\WarpTwistIesan\).
So: \[\mathbf{h}_{\varphi\psi} \neq -\Jsv\,\CenterShearLocation\]

The difference is precisely the \(\nabla\WarpTwistIesan\) coupling
terms. Define:
\[\bm{\lambda} := \int_\Section (\bm{r} - \CentroidLocation)(\bm{r} \cdot \nabla\WarpTwistIesan) \dA, \qquad \bm{\zeta} := \int_\Section r^2 \nabla\WarpTwistIesan \dA\]

Then:
\[\mathbf{h}_{\varphi\psi} = -\Jsv\,\CenterShearLocation + \frac{\nu}{2}\bm{\zeta} - \nu\bm{\lambda}\]
\end{Details}

The quantities \(\bm{\zeta}\) and \(\bm{\lambda}\) encode torsion-shear
coupling that exists even for symmetric sections when the origin is not
at the centroid. These do NOT vanish under the normalization
\(\int \WarpTwistIesan \dA = 0\).

\subsubsection{Derive X}

\textbf{Step 1: Compute \(\mathbf{L}_0^{-\top}\).}

We have: \[\mathbf{L}_0^{-1} = \begin{pmatrix}
1 & \mathbf{0}^\top & 0 & \mathbf{0}^\top \\
\mathbf{0} & \mathbf{0} & \mathbf{0} & \oneS \\
0 & -\CenterShearLocation^\top & 1 & \mathbf{0}^\top \\
\mathbf{0} & \oneS & \mathbf{0} & \mathbf{0}
\end{pmatrix}
\qquad
\mathbf{L}_0^{-\top} = \begin{pmatrix}
1 & \mathbf{0}^\top & 0 & \mathbf{0}^\top \\
\mathbf{0} & \mathbf{0} & -\CenterShearLocation & \oneS \\
0 & \mathbf{0}^\top & 1 & \mathbf{0}^\top \\
\mathbf{0} & \oneS & \mathbf{0} & \mathbf{0}
\end{pmatrix}\]

\textbf{Step 2: Compute \(\mathbf{R}^\top\).}

From the resultant matrix with partition
\(\bm{s} = (N, \bm{F}, T, \bm{M})\),
\(\bm{p} = (a_\varepsilon, \bm{a}_\Omega, a_\vartheta, \bm{b}_{\Section})\):

\[\mathbf{R}^\top = \begin{pmatrix}
EA & \mathbf{0}^\top & 0 & EA(\FrameBasis\times\CentroidLocation)^\top \\
EA\CentroidLocation & \mathbf{0} & \mathbf{0} & E(\FrameBasis^\times\IntRoR)^\top \\
0 & \mathbf{0}^\top & G\Jsv & \mathbf{0}^\top \\
\mathbf{0} & E\IntRGoRG & \mathbf{0} & \mathbf{0}
\end{pmatrix}\]

Using
\((\FrameBasis\times\CentroidLocation)^\top = -\CentroidLocation^\top\FrameBasis^\times\)
and \((\FrameBasis^\times\IntRoR)^\top = -\IntRoR\FrameBasis^\times\):

\[\mathbf{R}^\top = \begin{pmatrix}
EA & \mathbf{0}^\top & 0 & -EA\CentroidLocation^\top\FrameBasis^\times \\
EA\CentroidLocation & \mathbf{0} & \mathbf{0} & -E\IntRoR\FrameBasis^\times \\
0 & \mathbf{0}^\top & G\Jsv & \mathbf{0}^\top \\
\mathbf{0} & E\IntRGoRG & \mathbf{0} & \mathbf{0}
\end{pmatrix}
\]

\subparagraph{Step 3: Compute \(\mathbf{L}_0^{-\top}\mathbf{R}^\top\).}

\[
\mathbf{L}_0^{-\top}\mathbf{R}^\top = \begin{pmatrix}
EA & \mathbf{0}^\top & 0 & -EA\CentroidLocation^\top\FrameBasis^\times \\
\mathbf{0} & E\IntRGoRG & -G\Jsv\CenterShearLocation & \mathbf{0} \\
0 & \mathbf{0}^\top & G\Jsv & \mathbf{0}^\top \\
EA\CentroidLocation & \mathbf{0} & \mathbf{0} & -E\IntRoR\FrameBasis^\times
\end{pmatrix}
\]

\subparagraph{Step 4: Invert \(\bar{\mathbf{C}}\).}

The energy matrix has block structure:
\[\bar{\mathbf{C}} = \begin{pmatrix}
EA & \mathbf{0}^\top & 0 & EA\CentroidLocation^\top \\
\mathbf{0} & G\mathbf{H}_{\psi\psi} & G\mathbf{h}_{\varphi\psi} & \mathbf{0} \\
0 & G\mathbf{h}_{\varphi\psi}^\top & G\Jsv & \mathbf{0}^\top \\
EA\CentroidLocation & \mathbf{0} & \mathbf{0} & E\IntRoR
\end{pmatrix}\]

Reordering to \((\epsilon_0, \bm{\epsilon}_1, \bm{\beta}, \tau)\)
reveals block-diagonal structure:

\[
\bar{\mathbf{C}}_{\rm perm} = \begin{pmatrix}
\mathbf{A} & \mathbf{0} \\
\mathbf{0} & \mathbf{D}
\end{pmatrix}
\qquad\text{with:}\qquad
\mathbf{A} = \begin{pmatrix} EA & EA\CentroidLocation^\top \\ EA\CentroidLocation & E\IntRoR 
\end{pmatrix}, 
\quad 
\mathbf{D} = G\begin{pmatrix} \mathbf{H}_{\psi\psi} & \mathbf{h}_{\varphi\psi} \\ \mathbf{h}_{\varphi\psi}^\top & \Jsv \end{pmatrix}
\]

\textbf{Inverting \(\mathbf{A}\):} Schur complement is
\(E\IntRoR - EA\CentroidLocation\otimes\CentroidLocation = E\IntRGoRG\):
\[\mathbf{A}^{-1} = \begin{pmatrix} 
\frac{1}{EA} + \CentroidLocation^\top(E\IntRGoRG)^{-1}\CentroidLocation & -\CentroidLocation^\top(E\IntRGoRG)^{-1} \\
-(E\IntRGoRG)^{-1}\CentroidLocation & (E\IntRGoRG)^{-1}
\end{pmatrix},
\qquad\text{and}\qquad 
\mathbf{D}^{-1} = \frac{1}{G}\begin{pmatrix}
\mathbf{H}_{\psi\psi}^{-1} + \frac{1}{\tilde{\Jsv}}\bm{\eta}\otimes\bm{\eta} & -\frac{1}{\tilde{\Jsv}}\bm{\eta} \\
-\frac{1}{\tilde{\Jsv}}\bm{\eta}^\top & \frac{1}{\tilde{\Jsv}}
\end{pmatrix}\]

\[\tilde{\Jsv} := \Jsv - \mathbf{h}_{\varphi\psi}^\top\mathbf{H}_{\psi\psi}^{-1}\mathbf{h}_{\varphi\psi}
\qquad\text{and}\qquad 
\bm{\eta} := \mathbf{H}_{\psi\psi}^{-1}\mathbf{h}_{\varphi\psi}
\]


Permuting back: \[\bar{\mathbf{C}}^{-1} = \begin{pmatrix}
\frac{1}{EA} + \CentroidLocation^\top(E\IntRGoRG)^{-1}\CentroidLocation & \mathbf{0}^\top & 0 & -\CentroidLocation^\top(E\IntRGoRG)^{-1} \\
\mathbf{0} & \frac{1}{G}(\mathbf{H}_{\psi\psi}^{-1} + \frac{1}{\tilde{\Jsv}}\bm{\eta}\otimes\bm{\eta}) & -\frac{1}{G\tilde{\Jsv}}\bm{\eta} & \mathbf{0} \\
0 & -\frac{1}{G\tilde{\Jsv}}\bm{\eta}^\top & \frac{1}{G\tilde{\Jsv}} & \mathbf{0}^\top \\
-(E\IntRGoRG)^{-1}\CentroidLocation & \mathbf{0} & \mathbf{0} & (E\IntRGoRG)^{-1}
\end{pmatrix}\]

\textbf{Final result:}

\[
\boxed{\mathbf{X}^{-1} = \begin{pmatrix}
1 & \mathbf{0}^\top & 0 & \mathbf{0}^\top \\[4pt]
\mathbf{0} & \frac{E}{G}(\mathbf{H}_{\psi\psi}^{-1} + \frac{1}{\tilde{\Jsv}}\bm{\eta}\otimes\bm{\eta})\IntRGoRG & -\Jsv\mathbf{H}_{\psi\psi}^{-1}\CenterShearLocation - \frac{\Jsv}{\tilde{\Jsv}}\bm{\eta}(1 + \bm{\eta}^\top\CenterShearLocation) & \mathbf{0} \\[4pt]
0 & -\frac{E}{G\tilde{\Jsv}}\bm{\eta}^\top\IntRGoRG & \frac{\Jsv}{\tilde{\Jsv}}(1 + \bm{\eta}^\top\CenterShearLocation) & \mathbf{0}^\top \\[4pt]
\mathbf{0} & \mathbf{0} & \mathbf{0} & -\FrameBasis^\times
\end{pmatrix}}\]

where 
\(\bm{\eta}\) is
the shear-torsion coupling vector \eqref{eq:def-eta}, 
\(\tilde{\Jsv}\)
is the reduced torsion constant \eqref{eq}, and 
\(\mathbf{h}_{\varphi\psi}\) from \eqref{eq}.

\subsubsection{Structure of X}

\textbf{Local implementability:} Yes, automatically!

By construction, we defined \(\bm{e} = \mathbf{X}\hat{\bm{e}}\) where
\(\hat{\bm{e}}(x) = \mathbf{L}(x)\bm{p} = (\mathbf{L}_0 + x\mathbf{L}_1)\bm{p}\).
Therefore:
\[\TraceDeDpRoot = \mathbf{X}\mathbf{L}_0, \qquad \TraceDeDpOne = \mathbf{X}\mathbf{L}_1\]

The implementability condition
\(\TraceDeDpOne = \TraceDeDpRoot\mathbf{L}_0^{-1}\mathbf{L}_1\) is
satisfied identically:
\[\TraceDeDpRoot\mathbf{L}_0^{-1}\mathbf{L}_1 = \mathbf{X}\mathbf{L}_0\mathbf{L}_0^{-1}\mathbf{L}_1 = \mathbf{X}\mathbf{L}_1 = \TraceDeDpOne\]


\textbf{Plane-section property:} No!

To check this, we need \(\TraceDeDpRoot = \mathbf{X}\mathbf{L}_0\).
Looking at the \((5\text{-}6, 5\text{-}6)\) block of
\(\mathbf{X}^{-1}\):
\[[\mathbf{X}^{-1}]_{(5:6, 5:6)} = -\FrameBasis^\times\]

Inverting:
\([\mathbf{X}]_{(5:6, 5:6)} = (-\FrameBasis^\times)^{-1} = \FrameBasis^\times\)

Since \([\mathbf{L}_0]_{(5:6, 2:3)} = \oneS\) (the
\(\bm{a}_\Omega \mapsto \bm{\epsilon}_1\) block), we get:
\[[\TraceDeDpRoot]_{(5:6, 2:3)} = \FrameBasis^\times \cdot \oneS = \FrameBasis^\times\]

But the plane-section property requires
\([\TraceDeDpRoot]_{(5:6, 2:3)} = -\FrameBasis^\times\) (curvature from
bending should be
\(\bm{\kappa}_\perp = -\FrameBasis \times \bm{a}_\Omega\)).

\textbf{The sign is wrong.} The energy-conjugate trace does not satisfy
the plane-section property.

\begin{center}\rule{0.5\linewidth}{0.5pt}\end{center}


You're absolutely right to push on this. Let me think carefully about
what we actually have.

\textbf{The core issue:}

We derived \(\mathbf{X}^{-1}\) from the energy conjugacy condition at a
single section. But the resultants themselves depend on \(x\)!
Specifically, for a Saint-Venant solution:

\begin{itemize}
\tightlist
\item
  \(N, \bm{F}, T\) are constant along the beam
\item
  But \(\bm{M}(x) = \bm{M}_0 + x(\FrameBasis \times \bm{F})\) varies
  linearly
\end{itemize}

This is just equilibrium: \(\bm{M}' = \FrameBasis \times \bm{F}\).

\textbf{Let me redo the energy conjugacy condition properly:}

At section \(x\), energy conjugacy requires:
\[\bm{s}(x)^\top \delta\bm{e}(x) = \int_\Section \bm{\sigma}^\top \delta\bm{\varepsilon} \dA = \delta\hat{\bm{e}}(x)^\top \bar{\mathbf{C}} \hat{\bm{e}}(x)\]

If we assume \(\bm{e}(x) = \mathbf{X}\hat{\bm{e}}(x)\) with
\(\mathbf{X}\) independent of \(x\), then:
\[\mathbf{X}^\top \bm{s}(x) = \bar{\mathbf{C}} \hat{\bm{e}}(x)\]

Now express both sides in terms of \(\bm{p}\): -
\(\hat{\bm{e}}(x) = (\mathbf{L}_0 + x\mathbf{L}_1)\bm{p}\) -
\(\bm{s}(x) = -(\IesanStiff + x\mathbf{R}_1)\bm{p}\)

where \(\mathbf{R}_1\) captures the moment variation:
\[\mathbf{R}_1 = \begin{pmatrix}
0 & \mathbf{0}^\top & 0 & \mathbf{0}^\top \\
\mathbf{0} & \mathbf{0} & \mathbf{0} & \mathbf{0} \\
0 & \mathbf{0}^\top & 0 & \mathbf{0}^\top \\
\mathbf{0} & \mathbf{0} & \mathbf{0} & -\FrameBasis^\times E\IntRGoRG
\end{pmatrix}\]

The energy conjugacy condition becomes:
\[-\mathbf{X}^\top (\IesanStiff + x\mathbf{R}_1)\bm{p} = \bar{\mathbf{C}} (\mathbf{L}_0 + x\mathbf{L}_1)\bm{p}\]

Matching powers of \(x\): - \(x^0\):
\(-\mathbf{X}^\top \IesanStiff = \bar{\mathbf{C}} \mathbf{L}_0\) -
\(x^1\):
\(-\mathbf{X}^\top \mathbf{R}_1 = \bar{\mathbf{C}} \mathbf{L}_1\)

From \(x^0\):
\(\mathbf{X}^\top = -\bar{\mathbf{C}} \mathbf{L}_0 \IesanStiff^{-1}\)

Substituting into \(x^1\):
\[\bar{\mathbf{C}} \mathbf{L}_0 \IesanStiff^{-1} \mathbf{R}_1 = \bar{\mathbf{C}} \mathbf{L}_1\]

\textbf{This is a consistency condition.} If it holds, the
energy-conjugate trace with constant \(\mathbf{X}\) is both
implementable and exact. If not, we cannot have both.

\subsubsection{Consistency condition:}

We need: \(\mathbf{L}_0 \IesanStiff^{-1} \mathbf{R}_1 = \mathbf{L}_1\)
or equivalently:
\(\IesanStiff^{-1} \mathbf{R}_1 = \mathbf{L}_0^{-1} \mathbf{L}_1\)

We already computed: \[\mathbf{L}_0^{-1}\mathbf{L}_1 = \begin{pmatrix}
0 & \mathbf{0}^\top & 0 & -\CentroidLocation^\top \\
\mathbf{0} & \mathbf{0} & \mathbf{0} & \mathbf{0} \\
0 & \mathbf{0}^\top & 0 & \mathbf{0}^\top \\
\mathbf{0} & \mathbf{0} & \mathbf{0} & \oneS
\end{pmatrix}\]

Now I need \(\IesanStiff^{-1}\mathbf{R}_1\). Recall:
\[
\IesanStiff = \begin{pmatrix}
EA & EA\CentroidLocation^\top & 0 & \mathbf{0}^\top \\
\mathbf{0} & \mathbf{0} & \mathbf{0} & E\IntRGoRG \\
0 & \mathbf{0}^\top & G\Jsv & \mathbf{0}^\top \\
EA(\FrameBasis\times\CentroidLocation) & E\FrameBasis^\times\IntRoR & \mathbf{0} & \mathbf{0}
\end{pmatrix},
\qquad 
\mathbf{R}_1 = \begin{pmatrix}
0 & \mathbf{0}^\top & 0 & \mathbf{0}^\top \\
\mathbf{0} & \mathbf{0} & \mathbf{0} & \mathbf{0} \\
0 & \mathbf{0}^\top & 0 & \mathbf{0}^\top \\
\mathbf{0} & \mathbf{0} & \mathbf{0} & -\FrameBasis^\times E\IntRGoRG
\end{pmatrix}
\]


The relevant part of \(\IesanStiff^{-1}\) is
\([\IesanStiff^{-1}]_{(2,4)} = (E\IntRGoRG)^{-1}\).

So:
\[[\IesanStiff^{-1}\mathbf{R}_1]_{(2,4)} = (E\IntRGoRG)^{-1} \cdot (-\FrameBasis^\times E\IntRGoRG) = -\FrameBasis^\times\]

But \([\mathbf{L}_0^{-1}\mathbf{L}_1]_{(2,4)} = \mathbf{0}\).

\textbf{The consistency condition fails!}


The energy-conjugate trace as derived gives
\(\TraceDeDpOne \neq \TraceDeDpRoot \mathbf{L}_0^{-1}\mathbf{L}_1\), so it
has nonzero locality defect. Or if we force implementability by setting
\(\TraceDeDpOne = \TraceDeDpRoot \mathbf{L}_0^{-1}\mathbf{L}_1\), then it's
not exactly energy-conjugate at all \(x\).

Let me compute these systematically.

\subsubsection{Structure of X:}

From the block structure of \(\mathbf{X}^{-1}\):
\[\mathbf{X}^{-1} = \begin{pmatrix}
1 & \mathbf{0}^\top & 0 & \mathbf{0}^\top \\
\mathbf{0} & \mathbf{A} & \mathbf{b} & \mathbf{0} \\
0 & \mathbf{c}^\top & d & \mathbf{0}^\top \\
\mathbf{0} & \mathbf{0} & \mathbf{0} & -\FrameBasis^\times
\end{pmatrix},
\qquad\text{and}\qquad 
\mathbf{X} = \begin{pmatrix}
1 & \mathbf{0}^\top & 0 & \mathbf{0}^\top \\
\mathbf{0} & \tilde{\mathbf{A}} & \tilde{\mathbf{b}} & \mathbf{0} \\
0 & \tilde{\mathbf{c}}^\top & \tilde{d} & \mathbf{0}^\top \\
\mathbf{0} & \mathbf{0} & \mathbf{0} & \FrameBasis^\times
\end{pmatrix}
\]

where: 
\begin{itemize}
\item 
\(\mathbf{A} := \frac{E}{G}\Big(\mathbf{H}_{\psi\psi}^{-1} + \frac{\bm{\eta}\otimes\bm{\eta}}{\tilde{\Jsv}}\Big)\IntRGoRG\)
\item
\(\mathbf{b} := -\Jsv\mathbf{H}_{\psi\psi}^{-1}\CenterShearLocation - \frac{\Jsv}{\tilde{\Jsv}}\bm{\eta}(1 + \bm{\eta}^\top\CenterShearLocation)\)
\item \(\mathbf{c}^\top := \frac{E}{G\tilde{\Jsv}}\bm{\eta}^\top\IntRGoRG\) -
\(d := -\frac{\Jsv}{\tilde{\Jsv}}(1 + \bm{\eta}^\top\CenterShearLocation)\)
\end{itemize}


and
\(\begin{pmatrix} \tilde{\mathbf{A}} & \tilde{\mathbf{b}} \\ \tilde{\mathbf{c}}^\top & \tilde{d} \end{pmatrix} = \begin{pmatrix} \mathbf{A} & \mathbf{b} \\ \mathbf{c}^\top & d \end{pmatrix}^{-1}\).

Let me invert the central \(3 \times 3\) block systematically.

\textbf{The central block of \(\mathbf{X}^{-1}\):}
\[\begin{pmatrix} \mathbf{A} & \mathbf{b} \\ \mathbf{c}^\top & d \end{pmatrix}\]

with: -
\(\mathbf{A} = \frac{E}{G}\Big(\mathbf{H}_{\psi\psi}^{-1} + \frac{\bm{\eta}\bm{\eta}^\top}{\tilde{\Jsv}}\Big)\IntRGoRG\)
-
\(\mathbf{b} = -\Jsv\mathbf{H}_{\psi\psi}^{-1}\CenterShearLocation - \frac{\Jsv}{\tilde{\Jsv}}\bm{\eta}(1 + \bm{\eta}^\top\CenterShearLocation)\)
- \(\mathbf{c}^\top = \frac{E}{G\tilde{\Jsv}}\bm{\eta}^\top\IntRGoRG\) -
\(d = -\frac{\Jsv}{\tilde{\Jsv}}(1 + \bm{\eta}^\top\CenterShearLocation)\)

\textbf{Step 1: Compute \(\mathbf{A}^{-1}\).}

Using Sherman-Morrison on
\(\mathbf{H}_{\psi\psi}^{-1} + \frac{\bm{\eta}\bm{\eta}^\top}{\tilde{\Jsv}}\):

With
\(\bm{\eta}^\top\mathbf{H}_{\psi\psi}\bm{\eta} = \mathbf{h}_{\varphi\psi}^\top\mathbf{H}_{\psi\psi}^{-1}\mathbf{h}_{\varphi\psi} = \Jsv - \tilde{\Jsv}\):

\[\Big(\mathbf{H}_{\psi\psi}^{-1} + \frac{\bm{\eta}\bm{\eta}^\top}{\tilde{\Jsv}}\Big)^{-1} = \mathbf{H}_{\psi\psi} - \frac{\mathbf{H}_{\psi\psi}\bm{\eta}\bm{\eta}^\top\mathbf{H}_{\psi\psi}}{\tilde{\Jsv} + \Jsv - \tilde{\Jsv}} = \mathbf{H}_{\psi\psi} - \frac{\mathbf{h}_{\varphi\psi}\mathbf{h}_{\varphi\psi}^\top}{\Jsv}\]

Therefore:
\[\boxed{\mathbf{A}^{-1} = \frac{G}{E}\IntRGoRG^{-1}\Big(\mathbf{H}_{\psi\psi} - \frac{\mathbf{h}_{\varphi\psi}\mathbf{h}_{\varphi\psi}^\top}{\Jsv}\Big)}\]

\textbf{Step 2: Compute \(\mathbf{A}^{-1}\mathbf{b}\).}

Recall
\(\bm{\xi} := \mathbf{h}_{\varphi\psi} + \Jsv\CenterShearLocation = \frac{\nu}{2}\bm{\zeta} - \nu\bm{\lambda}\)
(the Poisson-torsion coupling).

After lengthy but straightforward algebra:
\[\Big(\mathbf{H}_{\psi\psi} - \frac{\mathbf{h}_{\varphi\psi}\mathbf{h}_{\varphi\psi}^\top}{\Jsv}\Big)\Big(\Jsv\mathbf{H}_{\psi\psi}^{-1}\CenterShearLocation + \frac{\Jsv\sigma}{\tilde{\Jsv}}\bm{\eta}\Big) = \Jsv\CenterShearLocation + \mathbf{h}_{\varphi\psi} = \bm{\xi}\]

where \(\sigma := 1 + \bm{\eta}^\top\CenterShearLocation\).

Therefore:
\[\boxed{\mathbf{A}^{-1}\mathbf{b} = -\frac{G}{E}\IntRGoRG^{-1}\bm{\xi}}\]

\textbf{Step 3: Compute \(\mathbf{c}^\top\mathbf{A}^{-1}\).}

Using
\(\bm{\eta}^\top\mathbf{H}_{\psi\psi} = \mathbf{h}_{\varphi\psi}^\top\)
and \(\bm{\eta}^\top\mathbf{h}_{\varphi\psi} = \Jsv - \tilde{\Jsv}\):

\[\mathbf{c}^\top\mathbf{A}^{-1} = \frac{1}{\tilde{\Jsv}}\bm{\eta}^\top\Big(\mathbf{H}_{\psi\psi} - \frac{\mathbf{h}_{\varphi\psi}\mathbf{h}_{\varphi\psi}^\top}{\Jsv}\Big) = \frac{1}{\tilde{\Jsv}}\Big(\mathbf{h}_{\varphi\psi}^\top - \frac{\Jsv - \tilde{\Jsv}}{\Jsv}\mathbf{h}_{\varphi\psi}^\top\Big)\]

Therefore:
\[\boxed{\mathbf{c}^\top\mathbf{A}^{-1} = \frac{\mathbf{h}_{\varphi\psi}^\top}{\Jsv}}\]

\textbf{Step 4: Compute the Schur complement
\(S = d - \mathbf{c}^\top\mathbf{A}^{-1}\mathbf{b}\).}

Define
\(\xi_{\Jsv} := \mathbf{h}_{\varphi\psi}^\top\IntRGoRG^{-1}\bm{\xi}\)  
(a
scalar coupling Poisson-torsion to bending stiffness).

\[\mathbf{c}^\top\mathbf{A}^{-1}\mathbf{b} = -\frac{G}{E\Jsv}\xi_{\Jsv}\]

\[S = -\frac{\Jsv\sigma}{\tilde{\Jsv}} + \frac{G\xi_{\Jsv}}{E\Jsv} = -\frac{E\Jsv^2\sigma - G\tilde{\Jsv}\xi_{\Jsv}}{E\Jsv\tilde{\Jsv}}\]

Define: \[\boxed{\Delta := E\Jsv^2\sigma - G\tilde{\Jsv}\xi_{\Jsv}}\]

Then \(S = -\frac{\Delta}{E\Jsv\tilde{\Jsv}}\) and:

\[
\begin{aligned}
\tilde{d} = S^{-1} &= -\frac{E\Jsv\tilde{\Jsv}}{\Delta}\\
\tilde{\mathbf{b}} &= -\mathbf{A}^{-1}\mathbf{b}\,\tilde{d} = -\frac{G\Jsv\tilde{\Jsv}}{\Delta}\IntRGoRG^{-1}\bm{\xi} \\
\tilde{\mathbf{c}}^\top &= -\tilde{d}\,\mathbf{c}^\top\mathbf{A}^{-1} = \frac{E\tilde{\Jsv}}{\Delta}\mathbf{h}_{\varphi\psi}^\top \\
\tilde{\mathbf{A}} &= \mathbf{A}^{-1} - \frac{\tilde{\mathbf{b}}\tilde{\mathbf{c}}^\top}{\tilde{d}} \\
&= \frac{G}{E}\IntRGoRG^{-1}\Big(\mathbf{H}_{\psi\psi} - \frac{\mathbf{h}_{\varphi\psi}\mathbf{h}_{\varphi\psi}^\top}{\Jsv} + \frac{G\tilde{\Jsv}}{\Delta}\bm{\xi}\mathbf{h}_{\varphi\psi}^\top\Big)
\end{aligned}
\]

\subsection{Trace Matrices}

\textbf{Computing \(\TraceDeDpRoot = \mathbf{X}\mathbf{L}_0\):}

With: \[\mathbf{L}_0 = \begin{pmatrix}
1 & \mathbf{0}^\top & 0 & \mathbf{0}^\top \\
\mathbf{0} & \mathbf{0} & \mathbf{0} & \oneS \\
0 & \mathbf{0}^\top & 1 & \CenterShearLocation^\top \\
\mathbf{0} & \oneS & \mathbf{0} & \mathbf{0}
\end{pmatrix},
\qquad
\mathbf{L}_1 = \begin{pmatrix}
0 & \mathbf{0}^\top & 0 & -\CentroidLocation^\top \\
\mathbf{0} & \mathbf{0} & \mathbf{0} & \mathbf{0} \\
0 & \mathbf{0}^\top & 0 & \mathbf{0}^\top \\
\mathbf{0} & \mathbf{0} & \mathbf{0} & \oneS
\end{pmatrix}
\]

Therefore: \[
\boxed{
\TraceDeDpRoot = \begin{pmatrix}
1 & \mathbf{0}^\top & 0 & \mathbf{0}^\top \\[4pt]
\mathbf{0} & \mathbf{0} & \tilde{\mathbf{b}} & \tilde{\mathbf{A}} + \tilde{\mathbf{b}}\CenterShearLocation^\top \\[4pt]
0 & \mathbf{0}^\top & \tilde{d} & \tilde{\mathbf{c}}^\top + \tilde{d}\CenterShearLocation^\top \\[4pt]
\mathbf{0} & \FrameBasis^\times & \mathbf{0} & \mathbf{0}
\end{pmatrix},
\qquad
\TraceDeDpOne = \begin{pmatrix}
0 & \mathbf{0}^\top & 0 & -\CentroidLocation^\top \\[4pt]
\mathbf{0} & \mathbf{0} & \mathbf{0} & \mathbf{0} \\[4pt]
0 & \mathbf{0}^\top & 0 & \mathbf{0}^\top \\[4pt]
\mathbf{0} & \mathbf{0} & \mathbf{0} & \FrameBasis^\times
\end{pmatrix}}
\]


\textbf{Physical interpretation:}

The fact that
it corresponds to a trace center at the origin is a \emph{consequence}
of the variational definition, not an input.

Specifically, the axial strain \(\epsilon\) is energy-conjugate
to the axial force \(N\), not to any moment. Pure bending
(\(\bm{a}_\Omega \neq 0\) with \(a_\varepsilon = 0\)) produces zero
axial force, so energy conjugacy requires
\(\delta\epsilon^{(a)} \cdot N = 0\), which is satisfied if
\(\epsilon^{(a)} = 0\).

\subsection{Trace Matrices (Corrected)}

Looking back at \([\mathbf{X}^{-1}]_{(5:6, 5:6)} = \FrameBasis^\times\),
we have:
\[[\mathbf{X}]_{(5:6, 5:6)} = (\FrameBasis^\times)^{-1} = -\FrameBasis^\times
\]


The corrected trace matrices are:

\[
\boxed{\TraceDeDpRoot = \begin{pmatrix}
1 & \mathbf{0}^\top & 0 & \mathbf{0}^\top \\[4pt]
\mathbf{0} & \mathbf{0} & \tilde{\mathbf{b}} & \tilde{\mathbf{A}} + \tilde{\mathbf{b}}\CenterShearLocation^\top \\[4pt]
0 & \mathbf{0}^\top & \tilde{d} & \tilde{\mathbf{c}}^\top + \tilde{d}\CenterShearLocation^\top \\[4pt]
\mathbf{0} & -\FrameBasis^\times & \mathbf{0} & \mathbf{0}
\end{pmatrix}
\qquad 
\TraceDeDpOne = \begin{pmatrix}
0 & \mathbf{0}^\top & 0 & -\CentroidLocation^\top \\[4pt]
\mathbf{0} & \mathbf{0} & \mathbf{0} & \mathbf{0} \\[4pt]
0 & \mathbf{0}^\top & 0 & \mathbf{0}^\top \\[4pt]
\mathbf{0} & \mathbf{0} & \mathbf{0} & -\FrameBasis^\times
\end{pmatrix}}\]

Now the \((4, 2:3)\) block of \(\TraceDeDpRoot\) is
\(-\FrameBasis^\times\), consistent with the plane-section property
\(\bm{\kappa}_\perp = -\FrameBasis \times \bm{a}_\Omega\).

\subsection{Correspondance}

Let me reconcile the two derivations carefully.

\subsubsection{The shear-flexure
relation}\label{the-shear-flexure-relation}

Your earlier derivation gives:
\[\bm{\gamma}^{\rm sh} = \BishearTensorRB^{-\top}\mathbf{H}\,\bm{b}_{\Section} \quad \Longleftrightarrow \quad \bm{b}_{\Section} = \mathbf{H}^{-1}\BishearTensorRB^\top\,\bm{\gamma}^{\rm sh}\]

From my \(\mathbf{X}^{-1}\), the \((2,2)\) block (ignoring torsion
coupling) gives:
\[\bm{\beta} = \frac{E}{G}\mathbf{H}_{\psi\psi}^{-1}\IntRGoRG\,\bm{\gamma}_\perp\]

\textbf{Reconciling via equilibrium:}

The key insight is that \(\BishearTensorRB\) (defined from the shear
stress distribution) equals \(E\IntRGoRG\) by equilibrium. Under pure
flexure loading \(\bm{b}_{\Section}\), the shear force is:
\[\bm{F} = \int_\Section G\,\bm{\varepsilon}_\gamma^{(b)} \dA = G\int_\Section \bm{\tau}_\psi \dA \cdot \bm{b}_{\Section} = \BishearTensorRB\,\bm{b}_{\Section}\]

But from the Ieşan equilibrium relation:
\[\bm{F} = E\IntRGoRG\,\bm{b}_{\Section}\]

Therefore: \[\boxed{\BishearTensorRB = E\IntRGoRG}\]

\textbf{Verifying consistency:}

Substituting into your formula:
\[\bm{b}_{\Section} = \mathbf{H}^{-1}\BishearTensorRB^\top\,\bm{\gamma}^{\rm sh} = (G\mathbf{H}_{\psi\psi})^{-1}(E\IntRGoRG)^\top\,\bm{\gamma}^{\rm sh} = \frac{E}{G}\mathbf{H}_{\psi\psi}^{-1}\IntRGoRG\,\bm{\gamma}^{\rm sh}\]

This matches my \([\mathbf{X}^{-1}]_{(2,2)}\) block exactly (when
torsion coupling vanishes). ✓

\textbf{The effective shear stiffness:}

Your formula:
\[\mathbf{K}_{\rm eff} = \BishearTensorRB\,\mathbf{H}^{-1}\BishearTensorRB^\top = (E\IntRGoRG)\,(G\mathbf{H}_{\psi\psi})^{-1}\,(E\IntRGoRG)^\top = \frac{E^2}{G}\IntRGoRG\,\mathbf{H}_{\psi\psi}^{-1}\,\IntRGoRG\]

And the shear force-strain relation:
\[\bm{F} = \mathbf{K}_{\rm eff}\,\bm{\gamma}^{\rm sh}\]

\textbf{Including torsion coupling:}

When \(\CenterShearLocation \neq \mathbf{0}\) or \(\bm{\eta} \neq \mathbf{0}\),
the full \((2,2)\) block of \(\mathbf{X}^{-1}\) is:
\[[\mathbf{X}^{-1}]_{(2,2)} = \frac{E}{G}\Big(\mathbf{H}_{\psi\psi}^{-1} + \frac{\bm{\eta}\bm{\eta}^\top}{\tilde{\Jsv}}\Big)\IntRGoRG\]


\subsection{Kinematic fields}

By integrating the strains.

\textbf{From curvature to rotation:}

We have \(\bm{\kappa} = \bm{\theta}'\), so:
\[\vartheta(x) = \vartheta(0) + \int_0^x \varkappa(s)\,ds\]
\[\bm{\theta}_\perp(x) = \bm{\theta}_\perp(0) + \int_0^x \bm{\kappa}_\perp(s)\,ds\]

From the trace matrices:
\[\varkappa = \tilde{d}\,a_\vartheta + (\tilde{\mathbf{c}}^\top + \tilde{d}\CenterShearLocation^\top)\bm{b}_{\Section} \quad \text{(constant in } x\text{)}\]
\[\bm{\kappa}_\perp(x) = -\FrameBasis^\times\bm{a}_\Omega - x\FrameBasis^\times\bm{b}_{\Section}\]

Integrating:
\[\vartheta(x) = \vartheta(0) + x\big(\tilde{d}\,a_\vartheta + (\tilde{\mathbf{c}}^\top + \tilde{d}\CenterShearLocation^\top)\bm{b}_{\Section}\big)\]
\[\bm{\theta}_\perp(x) = \bm{\theta}_\perp(0) - x\FrameBasis^\times\bm{a}_\Omega - \frac{1}{2}x^2\FrameBasis^\times\bm{b}_{\Section}\]

\textbf{From shear strain to displacement:}

We have \(\bm{\gamma} = \bm{u}' + \FrameBasis \times \bm{\theta}\).

For the axial part, \((\FrameBasis \times \bm{\theta})_\parallel = 0\),
so:
\[u_\parallel(x) = u_\parallel(0) + \int_0^x \epsilon(s)\,ds = u_\parallel(0) + a_\varepsilon x - \frac{1}{2}x^2\CentroidLocation^\top\bm{b}_{\Section}\]

For the transverse part,
\((\FrameBasis \times \bm{\theta})_\perp = \FrameBasis^\times\bm{\theta}_\perp\),
so:
\[\bm{u}_\perp' = \bm{\gamma}_\perp - \FrameBasis^\times\bm{\theta}_\perp\]

With
\(\bm{\gamma}_\perp = \tilde{\mathbf{b}}\,a_\vartheta + (\tilde{\mathbf{A}} + \tilde{\mathbf{b}}\CenterShearLocation^\top)\bm{b}_{\Section}\)
(constant):
\[\bm{u}_\perp(x) = \bm{u}_\perp(0) + x\bm{\gamma}_\perp - \FrameBasis^\times\int_0^x \bm{\theta}_\perp(s)\,ds\]

Computing the integral:
\[\int_0^x \bm{\theta}_\perp(s)\,ds = x\bm{\theta}_\perp(0) - \frac{1}{2}x^2\FrameBasis^\times\bm{a}_\Omega - \frac{1}{6}x^3\FrameBasis^\times\bm{b}_{\Section}\]

Using \((\FrameBasis^\times)^2 = -\mathbf{I}_2\):
\[\bm{u}_\perp(x) = \bm{u}_\perp(0) + x\bm{\gamma}_\perp - x\FrameBasis^\times\bm{\theta}_\perp(0) - \frac{1}{2}x^2\bm{a}_\Omega - \frac{1}{6}x^3\bm{b}_{\Section}\]

\textbf{Matching to the Ieşan solution:}

The section-rigid part of the Ieşan displacement at \(\bm{r} = \mathbf{0}\)
is:
\[\hat{\bm{u}}_{\rm rigid}(x, \mathbf{0}) = xa_\varepsilon\FrameBasis - \frac{1}{2}x^2(\CentroidLocation \cdot \bm{b}_{\Section})\FrameBasis - \frac{1}{2}x^2\bm{a}_\Omega - \frac{1}{6}x^3\bm{b}_{\Section}\]

The warping contribution at \(\bm{r} = \mathbf{0}\) is:
\[\hat{\bm{u}}_{\rm warp}(x, \mathbf{0}) = \WarpTwistIesan(\mathbf{0})(a_\vartheta + c_\vartheta)\FrameBasis + \bm{W}_{(a)}(\mathbf{0})\bm{a}_\Omega + \big(\bm{\Pi}_{(b)}(\mathbf{0}) + x\bm{\Pi}_{(b)}(\mathbf{0}) + \bm{\WarpShearIesan}(\mathbf{0})\big)\bm{b}_{\Section}\]

Wait---let me be more careful. The warping contributions depend on
whether the warping functions are normalized to vanish at the origin.

\textbf{Key observation:}

For the plane-section property with trace center at origin, we have:
\[\bm{u}(x) = \hat{\bm{u}}(x, \mathbf{0}) - \bm{\Phi}(\mathbf{0})\bm{\alpha}(x)\]

At \(x = 0\):
\[\bm{u}(0) = \hat{\bm{u}}(0, \mathbf{0}) - \bm{\Phi}(\mathbf{0})\bm{\alpha}(0)\]

The Ieşan displacement at \((0, \mathbf{0})\) includes: - Axial:
\(\WarpTwistIesan(\mathbf{0})\tau\,\FrameBasis\) (from torsion warping) -
Transverse:
\(\bm{W}_{(a)}(\mathbf{0})\bm{a}_\Omega + \bm{\Pi}_{(b)}(\mathbf{0})\bm{b}_{\Section} + \bm{\WarpShearIesan}(\mathbf{0})\bm{b}_{\Section}\)

\textbf{For a doubly-symmetric section with origin at centroid:}

All warping functions satisfy \(\WarpTwistIesan(\mathbf{0}) = 0\),
\(\bm{W}_{(a)}(\mathbf{0}) = \mathbf{0}\), \(\bm{\Pi}_{(b)}(\mathbf{0}) = \mathbf{0}\),
\(\bm{\WarpShearIesan}(\mathbf{0}) = \mathbf{0}\) by symmetry.

Then: \[\bm{u}(0) = \mathbf{0}, \qquad \bm{\theta}(0) = \mathbf{0}\]

The boundary condition \(\bm{u}(0) = \bm{\theta}(0) = \mathbf{0}\) means
``the origin is fixed and the section has no rotation''---a classical
clamped condition.

\textbf{For a general asymmetric section:}

The warping functions do not vanish at the origin. So:
\[\bm{u}(0) = \WarpTwistIesan(\mathbf{0})\tau\,\FrameBasis + \bm{W}_{(a)}(\mathbf{0})\bm{a}_\Omega + (\bm{\Pi}_{(b)}(\mathbf{0}) + \bm{\WarpShearIesan}(\mathbf{0}))\bm{b}_{\Section}\]

This is nonzero in general!

\subsubsection{Physical
interpretation:}\label{physical-interpretation-1}

The energy-conjugate trace with \(\bm{u}(0) = \bm{\theta}(0) = \mathbf{0}\)
does NOT mean ``the origin is fixed.'' It means:

\[\hat{\bm{u}}(0, \mathbf{0}) = \bm{\Phi}(\mathbf{0})\bm{\alpha}(0)\]

The origin can have nonzero displacement, provided it equals the warping
contribution there. This is an \textbf{implicit constraint coupling the
rigid and warping modes}.

Alternatively, if we want a true ``origin fixed'' boundary condition
\(\hat{\bm{u}}(0, \mathbf{0}) = \mathbf{0}\), then:
\[\bm{u}(0) = -\bm{\Phi}(\mathbf{0})\bm{\alpha}(0) \neq \mathbf{0}\]

\textbf{Summary:}

For asymmetric sections, the energy-conjugate trace boundary condition
\(\bm{u}(0) = \bm{\theta}(0) = \mathbf{0}\) does not correspond to fixing a
physical point. The trace center is at the coordinate origin, but the
warping evaluated there creates an offset between the beam-theoretic
``displacement'' and the actual 3D displacement of that point.

\subsection{Summary}


\textbf{Special case: \(\bm{\xi} = \mathbf{0}\) (doubly-symmetric section or
\(\nu = 0\)).}

When the Poisson-torsion coupling vanishes: -
\(\mathbf{h}_{\varphi\psi} = -\Jsv\CenterShearLocation\) -
\(\xi_{\Jsv} = 0\), so \(\Delta = E\Jsv^2\sigma\) -
\(\tilde{\mathbf{b}} = \mathbf{0}\) -
\(\tilde{d} = -\frac{\tilde{\Jsv}}{\Jsv\sigma}\) -
\(\tilde{\mathbf{c}}^\top = -\frac{\tilde{\Jsv}}{\Jsv\sigma}\CenterShearLocation^\top\)
-
\(\tilde{\mathbf{A}} = \frac{G}{E}\IntRGoRG^{-1}\Big(\mathbf{H}_{\psi\psi} - \Jsv\CenterShearLocation\CenterShearLocation^\top\Big)\)

\textbf{Further special case: \(\CenterShearLocation = \mathbf{0}\) (origin
at shear center).}

Then \(\sigma = 1\), \(\bm{\eta} = \mathbf{0}\), \(\tilde{\Jsv} = \Jsv\), and:
- \(\tilde{\mathbf{b}} = \mathbf{0}\) -
\(\tilde{\mathbf{c}}^\top = \mathbf{0}^\top\) - \(\tilde{d} = -1\) -
\(\tilde{\mathbf{A}} = \frac{G}{E}\IntRGoRG^{-1}\mathbf{H}_{\psi\psi}\)

The shear and torsion decouple completely.

% \clearpage 
% % Expressions for Paradiso's section locations $\bm{r}_*$ in the notation of draft.tex
\newcommand{\WarpTwistCentroid}{\bm{r}_{wtc}}
\newcommand{\ShearForceLocation}{\bm{\rho}_{F}}
\subsection*{Setup and Assumptions}

Work in centroidal coordinates so that the transverse origin is at the centroid:
\begin{equation}
  \CentroidLocation = \bar{\bm{\SectionLocation{}}} = \mathbf{0},
  \qquad
  \SectionCoordinate = \bm{r},
  \qquad
  \IntRGoRG = \mathbf{J}_{\!G} = \int_\Section \bm{r}\otimes\bm{r}\,\mathrm{d}A.
\end{equation}
If a non-centroidal frame is used, replace every occurrence of $\bm{r}$ by $\bm{r} - \CentroidLocation$.

All symbols such as $\AxialStress$, $\ShearStress$, $\ShearForce$, $\SectionCoordinate$, $\WarpTwistIesan$, $\WarpShearIesan$, $\BishearTensorPB$, $\IntRGoRG$, $\FrameBasis$, $\CenterShearLocation$, $\CenterTwistLocation$ etc.\ are those already defined in \texttt{draft.tex}.

\subsection*{1.\ $\bm{r}_F$ (center of normal stress)}

In Paradiso, for the axial stress field $\sigma_{\mathrm{ax}}$ and axial resultant $N$,
\begin{equation}
  \int_\Section \sigma_{\mathrm{ax}}\,\mathrm{d}A = N,
  \qquad
  \int_\Section \sigma_{\mathrm{ax}}\,\bm{r}\,\mathrm{d}A = N\,\bm{r}_F.
\end{equation}
In the notation of \texttt{draft.tex}, with $\AxialStress$ the axial stress and $\AxialRxn$ the axial resultant,
\begin{equation}
  \AxialRxn := \int_\Section \AxialStress\,\mathrm{d}A,
  \qquad
  \AxialRxn\,\bm{\SectionLocation}_F := \int_\Section \AxialStress\,\SectionCoordinate\,\mathrm{d}A.
\end{equation}
Thus
\begin{equation}
  \boxed{\bm{r}_F \;\longleftrightarrow\; \bm{\SectionLocation}_F
  \text{ defined by }
  \AxialRxn\,\bm{\SectionLocation}_F = \int_\Section \AxialStress\,\SectionCoordinate\,\mathrm{d}A.}
\end{equation}

\subsection*{2.\ $\bm{r}_P$ (prescribed load point)}

Paradiso defines $\bm{r}_P$ as the position in the cross-section where the shear force is applied. This is prescribed data, not determined by section integrals.

In \texttt{draft.tex}, the corresponding object is the \emph{trace center} $\bm{r}_{\mathcal T}$, i.e.\ the section location at which the trace line and FE node are placed:
\begin{equation}
  \boxed{\bm{r}_P \;\longleftrightarrow\; \bm{r}_{\mathcal T}.}
\end{equation}

\subsection*{3.\ $\bm{r}_p$ (statical equivalence point for the shear field)}

Paradiso defines $\bm{r}_p$ by the statical equivalence conditions for the shear stress field $\bm{\tau}_{\mathrm{sh}}$ associated with a shear force $\ShearForce$:
\begin{equation}
  \int_\Section \bm{\tau}_{\mathrm{sh}}\,\mathrm{d}A = \ShearForce,
\end{equation}
\begin{equation}
  \int_\Section \bm{r}\times\bm{\tau}_{\mathrm{sh}}\cdot\FrameBasis\,\mathrm{d}A
  = (\bm{r}_p\times\ShearForce)\cdot\FrameBasis.
\end{equation}

In \texttt{draft.tex} notation, with $\ShearStress$ the shear stress and $\ShearForce$ the shear resultant,
\begin{equation}
  \ShearForce := \int_\Section \ShearStress\,\mathrm{d}A,
\end{equation}
\begin{equation}
  (\bm{\SectionLocation}_p\times\ShearForce)\cdot\FrameBasis
  := \int_\Section \bigl(\SectionCoordinate\times\ShearStress\bigr)\cdot\FrameBasis\,\mathrm{d}A.
\end{equation}
Thus
\begin{equation}
  \boxed{\bm{r}_p \;\longleftrightarrow\; \bm{\SectionLocation}_p
  \text{ defined by }
  (\bm{\SectionLocation}_p\times\ShearForce)\cdot\FrameBasis
  = \int_\Section (\SectionCoordinate\times\ShearStress)\cdot\FrameBasis\,\mathrm{d}A.}
\end{equation}

\subsection*{4.\ $\bm{r}_T$ (center of twist)}

Paradiso defines the center of twist $\bm{r}_T$ by
\begin{equation}
  \bm{r}_T := \FrameBasis\times\mathbf{J}_G^{-1}\,\mathbf{s}_{\WarpTwistCentroid},
  \qquad
  \mathbf{s}_{\WarpTwistCentroid} := \int_\Section \WarpTwistCentroid\,\bm{r}\,\mathrm{d}A,
\end{equation}
where $\WarpTwistCentroid$ is the torsion warping scalar and $\mathbf{J}_G$ is the centroidal inertia tensor.

In \texttt{draft.tex}, with $\WarpTwistIesan$ the torsion warping in the Ieşan representation and $\IntRGoRG = \mathbf{J}_G$,
\begin{equation}
  \mathbf{s}_{\WarpTwistIesan} := \int_\Section \WarpTwistIesan\,\SectionCoordinate\,\mathrm{d}A,
\end{equation}
and the corresponding twist center is
\begin{equation}
  \boxed{
  \CenterTwistLocation
  = \FrameBasis\times\IntRGoRG^{-1}\,\mathbf{s}_{\WarpTwistIesan}
  = \FrameBasis\times\IntRGoRG^{-1}
    \int_\Section \WarpTwistIesan\,\SectionCoordinate\,\mathrm{d}A.}
\end{equation}
Thus
\begin{equation}
  \bm{r}_T \;\longleftrightarrow\; \CenterTwistLocation.
\end{equation}

\subsection*{5.\ $\bm{r}_\psi$ (shear/geometric center associated with shear warping)}

Paradiso defines a shear-related center $\bm{r}_\psi$ by
\begin{equation}
  \bm{r}_\psi
  = \FrameBasis\times\mathbf{J}_G^{-1}\,\bm{\gamma}_\psi,
  \qquad
  \bm{\gamma}_\psi
  = \int_\Section \mathbf{H}_\psi\bigl(\FrameBasis\times\bm{r}\bigr)\,\mathrm{d}A,
\end{equation}
where $\mathbf{H}_\psi := \bm{\psi}\otimes\nabla + \mathbf{A}^p$ is the shear warping operator built from the harmonic shear warping $\bm{\psi}$ and the isotropic particular solution $\mathbf{A}^p$.

In \texttt{draft.tex}, with
\begin{equation}
  \WarpShearIesan = \psi,
  \qquad
  \BishearTensorPB = \mathbf{H}_\psi,
\end{equation}
and again $\IntRGoRG = \mathbf{J}_G$, the corresponding center is
\begin{equation}
  \bm{\gamma}_\psi
  = \int_\Section \BishearTensorPB\bigl(\FrameBasis\times\SectionCoordinate\bigr)\,\mathrm{d}A,
\end{equation}
\begin{equation}
  \boxed{
  \CenterShearLocation
  = \FrameBasis\times\IntRGoRG^{-1}\,\bm{\gamma}_\psi
  = \FrameBasis\times\IntRGoRG^{-1}
    \int_\Section \BishearTensorPB\bigl(\FrameBasis\times\SectionCoordinate\bigr)\,\mathrm{d}A.}
\end{equation}
Thus
\begin{equation}
  \bm{r}_\psi \;\longleftrightarrow\; \CenterShearLocation.
\end{equation}

\subsection*{6.\ $\bm{r}_{\psi_p}$ (shear-warping shift relative to $\bm{r}_\psi$)}

Paradiso defines
\begin{equation}
  \bm{r}_{\psi_p} := \bm{r}_p - \bm{r}_\psi.
\end{equation}
Using the identifications above,
\begin{equation}
  \bm{r}_p \;\longleftrightarrow\; \bm{\SectionLocation}_p,
  \qquad
  \bm{r}_\psi \;\longleftrightarrow\; \CenterShearLocation,
\end{equation}
the corresponding vector in the notation of \texttt{draft.tex} is
\begin{equation}
  \boxed{
  \bm{\SectionLocation}_{\WarpShearIesan,p}
  := \bm{\SectionLocation}_p - \CenterShearLocation.}
\end{equation}
Thus
\begin{equation}
  \bm{r}_{\psi_p} \;\longleftrightarrow\; \bm{\SectionLocation}_p - \CenterShearLocation.
\end{equation}

\subsection*{7.\ Section locations not explicitly expressed}

The following objects from Paradiso are not expressed here purely in terms of the symbols of \texttt{draft.tex}, in order to avoid incomplete or speculative formulas:
\begin{itemize}
  \item The explicit integral expression for $\bm{\gamma}_\psi$ in terms of the section tensors
        $\mathbf{H}_{\psi\psi}$, $\mathbf{h}_{\varphi\psi}$, $J_0$ as they appear in Sec.~A.2.
  \item A closed-form expression for $\bm{r}_p$ solely in terms of those section tensors, obtained by
        substituting the explicit stress field into the statical equivalence conditions and solving.
\end{itemize}

